\documentclass[JEP,XML,SOM,Unicode,NoFloatCountersInSection,published]{cedram}
\datereceived{2024-05-16}
\dateaccepted{2024-07-18}
\dateepreuves{2024-07-22}

\usepackage[scr=rsfs,cal=euler]{mathalfa}
\renewcommand{\theenumi}{\roman{enumi}}
\hyphenation{dimen-sion depen-ding Essen-tial-ly}
\newcommand\sbw[2]{{{\mathsf{\Lambda}}^{{\!\scriptstyle#1}}{#2}}}

\usepackage{tikz,tkz-tab}

\DeclareMathOperator\codim{codim}
\DeclareMathOperator\Res{Res}

\newcommand\be{\boldsymbol{e}}
\newcommand\bx{\boldsymbol{x}}
\newcommand\by{\boldsymbol{y}}
\newcommand\bz{\boldsymbol{z}}
\newcommand\bX{\boldsymbol{X}}
\newcommand\bY{\boldsymbol{Y}}
\newcommand\bZZ{\boldsymbol{Z}}
\newcommand\bQ{\mathbb{Q}}
\newcommand\bN{\mathbb{N}}
\newcommand\bR{\mathbb{R}}
\newcommand\bZ{\mathbb{Z}}
\newcommand\cA{\mathcal{A}}
\newcommand\cB{\mathcal{B}}
\newcommand\cD{\mathfrak{D}}
\newcommand\cI{\mathcal{I}}
\newcommand\cU{\mathcal{U}}
\newcommand\cV{\mathcal{V}}
\newcommand\tP{\widetilde{P}}
\newcommand\tQ{\widetilde{Q}}
\newcommand\tR{\widetilde{R}}
\newcommand\tS{\widetilde{S}}
\newcommand\ee{\varepsilon}
\newcommand\GL{\mathrm{GL}}
\newcommand\GrO{\mathcal{O}}
\newcommand\hlambda{\widehat{\lambda}}
\newcommand\homega{\widehat{\omega}}\newcommand\HH{\mathcal{H}}
\newcommand\HHH{\mathcal{D}_\xi'}
\newcommand\HHstar{\mathcal{D}_{\xi}}
\newcommand{\norm}[1]{\left\|#1 \right\|_2}
\newcommand{\bnorm}[1]{\biggl\|#1 \biggr\|_2}
\newcommand{\normH}[1]{\|#1\|}
\newcommand{\psc}[2]{\left(#1 \mid #2\right)}
\newcommand{\Vect}[2][]{\left\la #2\right\ra_{#1}}

\let\subset \subseteq
\let\la\langle
\let\ra\rangle
\newcommand\AND{\quad\text{and}\quad}

\theoremstyle{definition}
\newtheorem{Def}{Definition}[section]
\newtheorem{Rem}[Def]{Remark}
\newtheorem*{remark*}{Remark}

\theoremstyle{plain}
\newtheorem{Prop}[Def]{Proposition}
\newtheorem{Lem}[Def]{Lemma}
\newtheorem{Thm}{Theorem}[section]
\newtheorem{Cor}[Def]{Corollary}

\datepublished{2024-07-23}
\begin{document}
\frontmatter

\title{On uniform polynomial approximation}

\author[\initial{A.} \lastname{Poëls}]{\firstname{Anthony} \lastname{Poëls}}
\address{Universite Claude Bernard Lyon 1, Institut Camille Jordan UMR 5208, \\
69622 Villeurbanne, France}
\email{poels@math.univ-lyon1.fr}
\urladdr{http://apoels-math-u.net/}

\subjclass{11J13, 11J82}

\keywords{Exponent of Diophantine approximation, heights, uniform polynomial approximation}

\begin{abstract}
Let $n$ be a positive integer and $\xi$ a transcendental real number. We are interested in bounding from above the uniform exponent of polynomial approximation $\widehat{\omega}_n(\xi)$. Davenport and Schmidt's original 1969 inequality $\widehat{\omega}_n(\xi)\leq 2n-1$ was improved recently, and the best upper bound known to date is $2n-2$ for each $n\geq 10$. In this paper, we develop new techniques leading us to the improved upper bound $2n-\frac{1}{3}n^{1/3}+\mathcal{O}(1)$.
\end{abstract}

\altkeywords{Exposants d'approximation diophantienne, hauteurs, approximation polynomiale uniforme}

\alttitle{Sur l'approximation polynomiale uniforme}

\begin{altabstract}
Soient $n$ un entier strictement positif et~$\xi$ un nombre réel transcendant. Nous cherchons à borner supérieurement l'exposant uniforme d'approximation polynomiale $\widehat{\omega}_n(\xi)$. Établie par Davenport et Schmidt en 1969, l'inégalité $\widehat{\omega}_n(\xi)\leq 2n-1$, a été améliorée pour la première fois récemment, et la meilleure borne supérieure connue à ce jour est $2n-2$ pour tout $n\geq 10$. Dans ce papier, nous développons de nouvelles techniques qui nous permettent d'obtenir la borne supérieure améliorée $2n-\frac{1}{3}n^{1/3}+\mathcal{O}(1)$.
\end{altabstract}

\maketitle
\tableofcontents

\mainmatter
\section{Introduction}

Let $\xi$ be a non-zero real number and let $n$ be a positive integer. Dirichlet's theorem (1842) is one of the most basic results of Diophantine approximation. It shows that for any real number $H > 1$, there exists a non-zero integer point $(x_0,\dots,x_n)\in\bZ^{n+1}$ such that
\begin{equation}
\label{eq: Dirichlet system}
\max\big\{|x_1|,\dots,|x_n|\big\} \leq H \AND |x_0 + x_1\xi + \cdots + x_n\xi^n| \leq H^{-n}.
\end{equation}
It is natural to ask if we can improve the exponent $n$ of $H^{-n}$, and this question gives rise to two Diophantine exponents. The so-called \emph{uniform exponent of approximation} $\homega_n(\xi)$ (\resp the \emph{ordinary} exponent $\omega_n(\xi)$), is the supremum of the real numbers $\omega > 0$ such that the system
\begin{equation*}
\normH{P} \leq H \AND 0 <|P(\xi)| \leq H^{-\omega}
\end{equation*}
admits a non-zero solution $P\in\bZ[X]$ of degree at most $n$ for each sufficiently large $H$ (\resp for arbitrarily large $H$). Here, $\normH{P}$ denotes the (naive) \emph{height} of $P$, defined as the largest absolute value of its coefficients. These quantities have been extensively studied over the past half-century, see for example \cite{bugeaud2015exponents} for a nice overview of the subject. By Dirichlet's theorem, if $\xi$ is not an algebraic number of degree $\leq n$, then we have
\[
\omega_n(\xi) \geq \homega_n(\xi) \geq n,
\]
and it is well known that those inequalities are equalities for almost all real numbers~$\xi$ (with respect to the Lebesgue measure). Note that if $\xi$ is an algebraic number of degree $d$, then $\homega_n(\xi)$ and $\omega_n(\xi)$ are both equal to $\min\{n,d-1\}$ (it is a consequence of Schmidt's subspace theorem, see \cite[Th.\,2.10]{bugeaud2015exponents}). We can therefore restrict our study to the set of transcendental real numbers. The initial question ``can we improve the exponent $n$ in Dirichlet's theorem?'' may be rephrased as follows: ``does there exist a transcendental real number $\xi$ satisfying $\homega_n(\xi) > n$?''. For $n=1$ the answer is negative and rather elementary to prove, so the first non-trivial case is $n=2$. Before the early $2000$s, it was conjectured that no such number existed. This belief was swept away by Roy's extremal numbers \cite{roy2003approximation}, \cite{roy2004approximation}, \cite{arbourRoyCriterionDegreTwo}, whose exponent $\homega_2$ is equal to the maximal possible value $(3+\sqrt 5)/2 = 2.618\cdots$. Since then, several families of transcendental real numbers whose uniform exponent $\homega_2$ is greater than $2$ have been discovered (see for example \cite{roy2007two}, \cite{bugeaud2005exponentsSturmian}, \cite{poels2017exponents, poelsExpoGeneralClass2021}). However, for $n \geq 3$ the mystery remains, and it is still an open question whether or not there exists $\xi\in\bR\setminus\overline{\bQ}$ with $\homega_n(\xi) > n$.

In this paper, we are interested in finding an upper bound for the uniform exponent $\homega_n(\xi)$, as this could provide clues to solving the initial problem. Brownawell's version of Gel'fond's criterion \cite{brownawell1974sequences} implies that $\homega_n(\xi) \leq 3n$. In 1969, Davenport and Schmidt \cite[Th.\,2b]{davenport1969approximation} showed that for any transcendental real number $\xi$ and any integer $n\geq 2$, we have
\begin{equation}
\label{eq: estimation Davenport and Schmidt}
\homega_n(\xi)\leq 2n-1.
\end{equation}
Up to now, few improvements have been made. Bugeaud and Schleischitz \cite[Th.\,2.1]{bugeaudSchleischitz2016Uniform} first got the upper bound
\begin{equation}
\label{eq: intro Bugeaud-Schleischitz}
\homega_n(\xi) \leq n-\frac{1}{2} + \sqrt{n^2-2n+1/4} = 2n-\frac{3}{2}+ \ee_n,
\end{equation}
where $\ee_n>0$ tends to $0$ as $n$ tends to infinity. Recently, Marnat and Moshchevitin \cite{marnat2018optimal} proved an important conjecture of Schmidt and Summerer on the ratio $\homega_n(\xi)/\omega_n(\xi)$ (see also \cite[Ch.\,2]{PhDMartin2019} for an alternative proof based on parametric geometry of numbers). In \cite{schleischitz2017some}, Schleischitz pointed out that we can use the aforementioned inequality in the proof of \eqref{eq: intro Bugeaud-Schleischitz} to get
\[
\homega_n(\xi) \leq 2n-2,
\]
for each $n\geq 10$. This is currently the best known upper bound. Let us also mention that using parametric geometry of numbers, Schleischitz \cite[Th.\,1.1]{schleischitz2017uniformPoly} was able to replace the estimate \eqref{eq: intro Bugeaud-Schleischitz} by
\begin{equation*}
\homega_n(\xi) \leq \frac{3(n-1)+\sqrt{n^2-2n+5}}{2} = 2n-2+ \ee_n'
\end{equation*}
where $\ee_n' > 0$ tends to $0$ as $n$ tends to infinity. For $n=3,\dots,9$, bounds that are better than \eqref{eq: estimation Davenport and Schmidt}, but nevertheless (strictly) greater than $2n-2$, are known. For example, it was proved in \cite{bugeaudSchleischitz2016Uniform} that for each transcendental real number $\xi$, we have
\[
\homega_3(\xi)\leq 3+\sqrt 2 = 4.41\cdots,
\]
see also the very recent work of Schleischitz \cite{schleischitz2024uniform}. In this paper, without relying on Marnat-Moshchevitin's inequality and with a different approach, we show in Section~\ref{Section: cas d = 2} that the upper bound $\homega_n(\xi)\leq 2n-2$ holds for any $n\geq 4$. We also improve the upper bound for~$\homega_3$.

\begin{Thm}
\label{Thm : main d=2}
Let $n\geq 3$ be an integer and $\xi\in\bR$ be a transcendental real number. If $n\geq 4$, then
\begin{equation*}
\homega_n(\xi) \leq 2n-2.
\end{equation*}
For $n=3$, we have the weaker estimate $\homega_3(\xi)\leq 2+\sqrt 5 = 4.23\cdots$.
\end{Thm}

We do not think that these upper bounds are optimal. It is interesting to note that Schleischitz, with a different method and under a technical condition, also found the estimates of Theorem~\ref{Thm : main d=2}, see~\cite{schleischitz2024uniform}. Our main result below is a significant improvement on the previous results as $n$ tends to infinity and does not require Marnat and Moshchevitin's inequality \cite{marnat2018optimal}.
\begin{Thm}
\label{Thm : main}
Set $a=1/3$. There exists a computable constant $N\geq 1$ such that, for each $n\geq N$ and any transcendental real number $\xi\in\bR$, we have
\begin{equation*}
\homega_n(\xi) \leq 2n- an^{1/3}.
\end{equation*}
\end{Thm}
The constant $a=1/3$ is not optimal. Numerical calculations based on the results from Section~\ref{section: proof of main thm} suggest that we could take $N$ rather ``small'' in Theorem~\ref{Thm : main} (maybe $N\leq 10^4$?). However, to keep the arguments and calculations as clear and simple as possible, we did not try to provide an explicit value of $N$.

Theorem~\ref{Thm : main} can be compared to \cite[Th.\,1.1]{poels2022simultaneous}, where we study $\hlambda_n(\xi)$, the uniform exponent of rational simultaneous approximation to the successive powers $\Xi=(1,\xi,\xi^2,\dots,\xi^n)$ (which is known to be, in a sense, dual to $\homega_n(\xi)$), see Section~\ref{section: notation} for the precise definition and more details. We were not able to deduce one result from the other, even though there are similarities in the arguments. For example, given a polynomial $P\in\bZ[X]$ of degree at most $n$, which is a good approximation, we can associate the $k+1$ polynomials $P,XP,\dots,X^{k}P$ of degree at most $n+k$. They provide information on $\homega_{n+k}(\xi)$. On the other hand, if we consider $\by\in\bZ^{n+1}$ which is a good approximation of $\Xi$ (for simultaneous approximation), we can associate the $k+1$ blocks of successive $n+1-k$ coordinates of $\by$, which are rather good approximations of $(1,\xi,\dots,\xi^{n-k})$. They in turn provide information on $\hlambda_{n-k}(\xi)$. Note that the difficulties in the proofs of both theorems are not in the same places. In particular, in this paper we have to work with \emph{irreducible} polynomials, a rather heavy constraint. Also, one of the most delicate parts of our approach is to bound from above the ordinary exponent $\omega_n(\xi)$, whereas this is rather ``simple'' to do for the ordinary exponent $\lambda_n(\xi)$ in \cite{poels2022simultaneous}.

Before presenting our strategy, let us quickly explain Davenport and Schmidt's proof of the upper bound \eqref{eq: estimation Davenport and Schmidt}. Given a real number $\homega < \homega_n(\xi)$, they show, using elementary means and Gelfond's lemma, that there are infinitely many pairs of coprime polynomials $P,Q\in \bZ[X]$ of degree at most $n$, such that
\[
\normH{Q} \leq \normH{P} \AND \max\{|Q(\xi)|, |P(\xi)|\} \ll \normH{P}^{-\homega},
\]
(where the implicit constant only depends on $n$). It implies that the resultant $\Res(P,Q)$, which is a non-zero integer, satisfies
\[
1 \leq |\Res(P,Q)| \ll \normH{P}^{n-1}\normH{Q}^{n-1} \max\big\{\normH{P}|Q(\xi)|, \normH{Q}|P(\xi)| \big\} \ll \normH{P}^{2n-1-\homega}.
\]
The first upper bound for $|\Res(P,Q)|$ is classical, see Lemma~\ref{lem: estimation classique du resultant}. Since $\normH{P}$ can be arbitrarily large, they deduced that the exponent $2n-1-\homega$ is non-negative. Estimate \eqref{eq: estimation Davenport and Schmidt} follows by letting $\homega$ tend to $\homega_n(\xi)$. Note that the term $2n$ in \eqref{eq: estimation Davenport and Schmidt} is directly related to the size of the $2n\times2n$ determinant defining $\Res(P,Q)$ (if we suppose that~$P$ and~$Q$ have degree exactly $n$).

The key idea in the proof of our main Theorem~\ref{Thm : main} is to work with a large number of ``good'' linearly independent polynomial approximations $Q_0, \dots Q_{j+1}$ rather than just two polynomials $P$ and $Q$ as above. By doing this, we can replace $\Res(P,Q)$ by a non-zero $(2n-j)\times (2n-j)$ determinant depending on the coefficients of $Q_0,\dots,Q_{j+1}$. Under the ideal and unlikely assumption that
\begin{equation}
\label{eq intro: controle des Q_i}
\normH{Q_k} \leq \normH{Q_0} \AND |Q_k(\xi)| \ll \normH{Q_0}^{-\homega} \qquad \textrm{(for $k=0,\dots,j$)},
\end{equation}
the aforementioned determinant would be bounded from above by $\normH{Q_0}^{2n-j-1 - \homega}$. So, together with an additional non-vanishing assumption, it would lead to $\homega_n(\xi) \leq 2n-j-1$. Several new difficulties arise when trying to make the above arguments work. We introduce the tools for the construction of the generalized resultant in Section~\ref{section: espaces V_N}. To ensure that this determinant does not vanish, we need the extra assumption that $Q_0,\dots,Q_{j+1}$ are irreducible polynomials. The idea is to first fix a sequence of best approximations, that we call \emph{minimal polynomials}, and then to consider their highest-degree irreducible factors (which also happen to be rather good approximations). We~deal with this question in Section~\ref{section: sequence des Q_i}. Two obstacles remain. Firstly, note that it may be possible that the best polynomial approximations span a subspace of dimension~$3$, even when $\xi$ is transcendental and $n$ is large, see \cite[Th.\,1.3]{moshchevitin2007best}. Therefore, as soon as $j > 1$ (we will later choose $j\asymp n^{1/3}$), we have to justify that we can find $j+2$ linearly independent polynomials as above. The second major problem is the control of the sequence $Q_0,\dots,Q_{j+1}$. Estimates \eqref{eq intro: controle des Q_i} seem out of reach, instead we get upper bounds of the form
\begin{equation}
\label{eq intro: controle des Q_i V2}
\normH{Q_k} \leq \normH{Q_0} \AND |Q_k(\xi)| \ll \normH{Q_0}|^{-\homega\theta} \qquad \textrm{(for $k=0,\dots,j$)},
\end{equation}
where $\theta < 1$ depends only on $n$ and $j$, and is ``close'' to $1$ if $j$ is not too large compared to $n$. The main ingredients for showing this are related to \emph{twisted heights}, see Sections~\ref{def petitesse HHstar(V)} and the appendix, and an important inequality on the height of subspaces due to Schmidt. The parameter $\theta$ in \eqref{eq intro: controle des Q_i V2} is a function of the exponent of best approximation $\omega_n(\xi)$. We show in Section~\ref{section: estimation de omega_n} that if the uniform exponent satisfies $\homega_n(\xi) \geq 2n-d$ (with $d \ll n^{1/3}$), then the ordinary exponent $\omega_n(\xi)$ is bounded from above by $2n+2d^2$, and the ratio $\homega_n(\xi)/\omega_n(\xi)$ is therefore close to~$1$. This part, which is essentially independent from the others, is rather delicate, because we work with the polynomials $Q_i$. They are certainly irreducible, but not as good approximations as the minimal polynomials. More precisely, there could be large gaps between the height of two successive $Q_i$. If we could drop the irreducibility condition and directly work with the sequence of minimal polynomials, we could possibly replace the upper bound $2n- \GrO(n^{1/3})$ with $2n- \GrO(n^{1/2})$ in Theorem~\ref{Thm : main}. Section~\ref{section: proof of main thm} is devoted to the proof of Theorem~\ref{Thm : main}.

\subsubsection*{Acknowledgements}
I would like to thank Damien Roy for his attentive reading of this work and his many comments, which helped to improve its overall presentation and clarity. I also thank the referee for their work and valuable comments.

\section{Notation}
\label{section: notation}
Throughout this paper, $\xi$ denotes a transcendental real number. The floor (\resp ceiling) function is denoted by $\lfloor \cdot \rfloor$ (\resp $\lceil \cdot \rceil$). If $f, g : I \to [0, +\infty)$ are two functions on a set $I$, we write $f = \GrO(g)$ or $f \ll g$ or $g \gg f$ to mean that there is a positive constant $c$ such that $f (x) \leq cg(x)$ for each $x\in I$. We write $f \asymp g$ when both $f \ll g$ and $g \ll f$ hold.

Let $K$ be a field. If $\cA$ is a subset of a $K$-vector space $V$, we denote by $\Vect[K]{\cA} \subset V$ the $K$-vector space spanned by $\cA$, with the convention that $\Vect[K]{\emptyset} = \{0\}$.

Given a ring $A$ (typically $A=\bR$ or $\bZ$) and an integer $n\geq 0$, we denote by $A[X]$ the ring of polynomials in $X$ with coefficients in $A$, and by $A[X]_{\leq n} \subset A[X]$ the subgroup of polynomials of degree at most $n$. We say that $P\in\bZ[X]$ is \emph{primitive} if it non-zero and the greatest common divisor of its coefficients is $1$. Given $P=\sum_{k=0}^{n} a_kX^k \in \bR[X]$, we set
\begin{equation*}
\normH{P} = \max_{0\leq k \leq n} |a_k|.
\end{equation*}

Gelfond's lemma is the following statement (see \eg \cite[Lem.\,A.3]{bugeaud2004approximation} as well as \cite{brownawell1974sequences}). For any non-zero polynomials $P_1,\dots,P_r\in\bR[X]$ with product $P = P_1\cdots P_r$ of degree at most $n$, we have
\begin{equation}
\label{eq: Gelfond's lemma}
e^{-n} \normH{P_1}\cdots \normH{P_r} < \normH{P} < e^{n} \normH{P_1}\cdots \normH{P_r}.
\end{equation}
In particular, for each non-zero polynomial $P \in \bZ[X]_{\leq n}$ and each factor $Q\in\bZ[X]$ of $P$, we have $e^{-n}\normH{Q} < \normH{P}$. We will often use \eqref{eq: Gelfond's lemma} as follows. If $Q\in\bZ[X]_{\leq n}$ is irreducible and if $P\in\bZ[X]_{\leq n}$ is a non-zero polynomial which satisfies $\normH{P} \leq e^{-n}\normH{Q}$, then $Q$ cannot divide $P$. They are thus coprime polynomials.

We recall the definition of the resultant, which, as explained in the introduction, is useful for estimating the exponent $\homega_n(\xi)$ (see also Section~\ref{section: resultant et premier resultat}). Let $P,Q\in\bZ[X]$ be non-constant polynomials of degree $p$ and $q$ respectively, and let $a_i,b_j\in \bZ$ such that $P(X) = \sum_{k=0}^{p} a_kX^k$ and $Q(X) = \sum_{k=0}^{q} b_kX^k$.
Their \emph{resultant} $\Res(P,Q)$ is defined as the $(q+p)$-dimensional determinant
\begin{equation}
\label{eq det resultant}
\arraycolsep3.5pt
\Res(P,Q) = \begin{array}{cc}
\left | \begin{array}{cccc}
a_p & 0 & \dots \\
a_{p-1} & a_p & \\
\vdots & \vdots & \ddots \\
a_0 & \\
0 & a_0 \\
\vdots & \vdots & \ddots \\
& & & a_0
\end{array}\right. &
\left. \hspace*{-10pt}
\begin{array}{cccc}
b_q & 0 & \dots \\
b_{q-1} & b_q & \\
\vdots & \vdots & \ddots \\
b_0 & \\
0 & b_0 \\
\vdots & \vdots & \ddots \\
& & & b_p
\end{array}\right| \\
\underbrace{\hspace*{2.7cm}}_{q} &
\underbrace{\hspace*{2.3cm}}_{p}
\end{array}.
\end{equation}
Besides the exponents of linear approximation $\omega_n$ and $\homega_n$, we will also need the following exponents of simultaneous rational approximation. For each positive integer~$n$, the exponent $\hlambda_n(\xi)$ (\resp $\lambda_n(\xi)$) is the supremum of the real numbers $\lambda \geq 0$ such that the system
\begin{equation*}
|y_0| \leq Y \AND L(\by) \leq Y^{-\lambda} \quad \textrm{where } L(\by):=\max_{1\le k \le n} |y_0\xi^k - y_k|,
\end{equation*}
admits a non-zero integer solution $\by=(y_0,\dots,y_n)\in\bZ^{n+1}$ for each sufficiently large $Y\geq 1$ (\resp for arbitrarily large $Y$). Dirichlet's theorem \cite[\S II.1, Th.\,1A]{schmidt1996diophantine} implies that $\hlambda_n(\xi)\geq 1/n$. The best upper bounds known to date for $\hlambda_n(\xi)$ when $n\geq 4$ are established in joint work with Roy in \cite{poels2022simultaneous}. In particular, there is an explicit positive constant $a$ such that
\[
\hlambda_n(\xi) \leq \frac{1}{n/2+an^{1/2}+1/3},
\]
and sharper results are also obtained when $n$ is small.

\section{Minimal polynomials}
\label{section: minimal pol}

A \emph{sequence of minimal polynomials} (associated to $n$ and $\xi$) is a sequence $(P_i)_{i\geq 0}$ of non-zero polynomials in $\bZ[X]_{\leq n}$ satisfying the following properties:
\begin{enumerate}
\item the sequence $\big(\normH{P_i}\big)_{i\geq 0}$ is strictly increasing;
\item the sequence $\big(|P_i(\xi)|\big)_{i\geq 0}$ is strictly decreasing;
\item if $|P(\xi)| < |P_i(\xi)|$ for some index $i\geq 0$ and a non-zero $P\in\bZ[X]_{\leq n}$, then $\normH{P}\geq \normH{P_{i+1}}$.
\end{enumerate}

Note that if we require the dominant coefficient of $P_i$ to be positive (and since $\xi$ is transcendental), then the above sequence is unique up to its first terms. Let $(P_i)_{i\geq 0}$ be a sequence as above. We have the classical formulas:
\begin{equation}
\label{eq: exposant via minimal points}
\homega_n(\xi) = \liminf_{i\to \infty} \frac{-\log |P_i(\xi)|}{\log \normH{P_{i+1}}} \AND \omega_n(\xi) = \limsup_{i\to \infty} \frac{-\log |P_i(\xi)|}{\log \normH{P_{i}}}.
\end{equation}
In particular, given a positive real number $\homega$ with $\homega < \homega_n(\xi)$, then we have, for each sufficiently large index $i$,
\begin{equation}
\label{eq: H(P_i+1) controle par H(P_i)}
|P_i(\xi)| \leq \normH{P_{i+1}}^{-\homega} \AND \normH{P_{i+1}}^\tau \leq \normH{P_i},\quad \textrm{where } \tau:= \frac{\homega}{\omega_n(\xi)},
\end{equation}
(with the convention $\tau = 0$ if $\omega_n(\xi) = \infty$). The second inequality in \eqref{eq: H(P_i+1) controle par H(P_i)} asks for an upper bound on $\omega_n(\xi)$. Given a non-zero $P \in \bZ[X]$, we set $\omega(P)=0$ if $\normH{P} = 1$. Otherwise, we denote by $\omega(P)$ the real number satisfying
\[
|P(\xi)| = \normH{P}^{-\omega(P)}.
\]
With this notation, we have
\begin{equation}
\label{eq: encadrement omega(P_i)}
\omega_n(\xi) = \limsup_{\substack{\normH{P}\to \infty \\ P\in\bZ[X]_{\leq n}}} \omega(P) = \limsup_{i\to \infty} \omega(P_i) \AND \liminf_{i\to \infty} \omega(P_i) \geq \homega_n(\xi).
\end{equation}
The following results are well-known. We prove them for the sake of completion. The first one follows from the arguments of the proof of \cite[Lem.\,2]{davenport1967approximation} (see also \cite[Lem.\,4.1]{roy2004approximation}).

\begin{Lem}
\label{Lem: P_i, P_i+1 est une Z-base}
Let $i\geq 0$ and write $V_i = \Vect[\bR]{P_i,P_{i+1}} \subset \bR[X]_{\leq n}$. Then $\{P_i, P_{i+1}\}$ forms a $\bZ$-basis of the lattice $V_i\cap \bZ[X]_{\leq n}$.
\end{Lem}

\begin{proof}
By contradiction, suppose that $\{P_i, P_{i+1}\}$ is not a $\bZ$-basis of $V_i\cap \bZ[X]_{\leq n}$. Then there exists a non-zero $Q\in \bZ[X]_{\leq n}$ which may be written as $Q = rP_i + sP_{i+1}$, where $r,s\in\bQ$ satisfy $|r|,|s|\leq 1/2$. In particular, we have
\begin{align*}
\normH{Q} &\leq |r|\normH{P_i} + |s|\normH{P_{i+1}} < \normH{P_{i+1}}, \\
|Q(\xi)| &\leq |r||P_i(\xi)| + |s||P_{i+1}(\xi)| < |P_{i}(\xi)|.
\end{align*}
This contradicts the minimality property of $P_i$.
\end{proof}

The next result is analogous to the second part of \cite[Lem.\,4.1]{roy2004approximation}. The construction of $S_i$ is due to Davenport and Schmidt \cite{davenport1967approximation}.

\begin{Lem}
\label{Lem: P_iH_i+1 = P_j-1H_j}
For each $i\geq 0$, define
\[
S_i = P_i(\xi)P_{i+1} - P_{i+1}(\xi)P_i \in\bR[X]_{\leq n}.
\]
Then
\[
\frac{1}{2}\normH{S_i} \leq \normH{P_{i+1}}|P_i(\xi)| \leq 2\normH{S_i}.
\]
Moreover, if for integers $0\leq i < j$ the space spanned by $P_i, P_{i+1}, \cdots, P_{j}$ has dimension~$2$, then $S_{j-1} = \pm S_i$. In particular
\[
\normH{P_{i+1}}|P_i(\xi)| \asymp \normH{P_{j}}|P_{j-1}(\xi)|.
\]
\end{Lem}

\begin{Rem}
Note that the quantity $\normH{S_i}$ satisfies $\normH{S_i} \asymp \HHstar(V_i)$, where $\HHstar$ is defined in Section~\ref{def petitesse HHstar(V)} and $V_i=\Vect[\bR]{P_i,P_{i+1}}$. We will study the function $\HHstar$ more deeply later.
\end{Rem}

\begin{proof}
We easily get $\normH{S_i} \leq 2\normH{P_{i+1}}|P_i(\xi)|$. Define $R_+, R_- \in\bZ[X]_{\leq n}$ by
\[
R_\pm = P_{i+1} \pm P_i.
\]
Suppose that there exists $\ee\in\{+,-\}$ such that $|R_\ee(\xi)| \leq |P_i(\xi)|/2$. Then by minimality of $P_i$, we must have $\normH{R_\ee} \geq \normH{P_{i+1}}$. Since $S_i = P_i(\xi)R_\ee - R_\ee(\xi) P_i$, we find
\begin{equation*}
\normH{S_i} \geq |P_i(\xi)|\normH{R_\ee} - |R_\ee(\xi)|\normH{P_i} \geq \frac{1}{2} \normH{P_{i+1}}|P_i(\xi)|.
\end{equation*}
Assume that $|R_+(\xi)|, |R_-(\xi)| \geq |P_i(\xi)|/2$. This is equivalent to
\[
|P_{i+1}(\xi)| \leq \frac{1}{2}|P_i(\xi)|.
\]
Again, this yields $\normH{S_i} \geq |P_i(\xi)|\normH{P_{i+1}} - |P_{i+1}(\xi)|\normH{P_i} \geq \normH{P_{i+1}}|P_i(\xi)|/2$.

Now, let us write $V_i = \Vect[\bR]{P_i,\dots,P_j}$, with $j > i$, and suppose that $V_i$ has dimension~$2$. We need to prove that $S_{j-1} = \pm S_i$. If $j=i+1$ it is automatic, we may therefore assume that $j\geq i+2$. By Lemma~\ref{Lem: P_i, P_i+1 est une Z-base}, there exist $a,b\in\bZ$ such that $P_i = aP_{i+1} + bP_{i+2}$. Since $\{P_i,P_{i+1}\}$ is also a $\bZ$-basis of $V_i$, we have $b=\pm 1$, and we deduce that
\[
S_i = \big(aP_{i+1}(\xi) + bP_{i+2}(\xi)\big)P_{i+1} - P_{i+1}(\xi)\big(aP_{i+1} + bP_{i+2}\big) = -bS_{i+1} = \pm S_{i+1}.
\]
By induction, we get $S_i = \pm S_{i+1} = \cdots = \pm S_{j-1}$.
\end{proof}

The proof of \cite[Lem.\,3]{davenport1967approximation} (which deals with the case $n=2$) yields the classical following result.

\begin{Lem}
\label{lem: I infini}
Suppose $n\geq 2$. Then, there are infinitely many indices $i\geq 1$ for which $P_{i-1}$, $P_i$ and $P_{i+1}$ are linearly independent.
\end{Lem}

\begin{proof}
By contradiction, suppose that there exists $i\geq 0$ such that $V=\Vect[\bR]{P_i,P_{i+1},\dots}$ has dimension $2$. By Lemma~\ref{Lem: P_iH_i+1 = P_j-1H_j} there exists $c>0$ such that for each $j > i$ we have
\[
0 < \normH{P_{i+1}}|P_i(\xi)| \leq c\normH{P_j}|P_{j-1}|.
\]
This leads to a contradiction since $\normH{P_j}|P_{j-1}| \leq \normH{P_j}^{1-\homega_n(\xi)+o(1)}$ tends to $0$ as $j$ tends to infinity.
\end{proof}

\begin{Rem}
As mentioned in the introduction, it is however possible that all polynomials $P_i$ with $i$ large enough lie in a subspace of dimension $3$ , see \cite[Th.\,1.3]{moshchevitin2007best}.
\end{Rem}

\section{Resultant and first estimates}
\label{section: resultant et premier resultat}

The following useful result can easily be derived from the proof of \cite[\S 5]{davenport1969approximation} (see also of \cite[Lem.\,1]{brownawell1974sequences}). We recall the arguments since they illustrate (in a simpler situation) how we will deal with generalized determinants.

\begin{Lem}
\label{lem: estimation classique du resultant}
Let $p,q$ be positive integers with $p,q\leq n$. There exists a constant $c>0$ depending on $\xi$ and $n$ only, with the following property. For any polynomials $P,Q\in\bZ[X]$ of degree $p$ and $q$ respectively, we have
\begin{equation*}
|\Res(P,Q)| \leq c \normH{P}^{q-1}\normH{Q}^{p-1} \max\big\{\normH{P}|Q(\xi)|, \normH{Q}|P(\xi)| \big\}.
\end{equation*}
\end{Lem}

\begin{proof}
Let $a_i,b_j\in \bZ$ such that $P(X) = \sum_{k=0}^{p} a_kX^k$ and $Q(X) = \sum_{k=0}^{q} b_kX^k$. For $i=1,\dots,p+q-1$, we add to the last row of the determinant \eqref{eq det resultant} its $i$-th row multiplied by $\xi^{p+q-i}$. This last row now becomes
\begin{equation*}
\Big(\xi^{q-1}P(\xi), \dots, \xi P(\xi), P(\xi), \xi^{p-1}Q(\xi), \dots, \xi Q(\xi), Q(\xi)\Big).
\end{equation*}
Using the upper bounds $|a_i| \leq \normH{P}$ and $|b_j|\leq \normH{Q}$ for the other entries of \eqref{eq det resultant}, we~obtain
\begin{equation*}
|\Res(P,Q)| \ll \normH{P}^{q-1}|P(\xi)| \normH{Q}^p + \normH{P}^q\normH{Q}^{p-1}|Q(\xi)|,
\end{equation*}
where the implicit constant only depends on $p, q$ and $\xi$.
\end{proof}

The next result, which is also based on inequalities involving resultants, will be used in Section~\ref{section: estimation de omega_n}. It ensures that if $R\in\bZ[X]$ is a ``good'' approximation, in the sense that $R(\xi)$ is very small compared to $\normH{R}$, and if we write $R$ as a product of coprime polynomials $B_1\cdots B_k$, then one of those factors is also a ``good'' approximation, while the product of the others is not.

\begin{Lem}
\label{lem: lem general 2}
Let $m,k$ be positive integers. There exists a constant $c>0$ depending on $m$ and $\xi$ only, with the following property. Let $B_1,\dots,B_k\in\bZ[X]$ be non-constant, pairwise coprime polynomials, and suppose that $R:= B_1\cdots B_k$ has degree at most $m$. Then, there exists $j\in\{1,\dots,k\}$ such that
\[
|B_j(\xi)| \leq c\normH{R}^{m-1}|R(\xi)| \AND \prod_{\substack{i=1 \\ i\neq j}}^{k} |B_i(\xi)| \geq c^{-1}\normH{R}^{-(m-1)}.
\]
\end{Lem}

\begin{proof}
If $k=1$ this is trivial. We now suppose that $k\geq 2$ and write $d_i = \deg(B_i)$ for $i=1,\dots, k$. By hypothesis, we have $\deg(R) = d_1+\cdots + d_k \leq m$. Note that
\begin{equation}
\label{eq proof: log H(P) = sum log H(R_i)}
|R(\xi)| = \prod_{i=1}^{k} |B_i(\xi)| \AND \normH{R} \asymp \prod_{i=1}^{k} \normH{B_i},
\end{equation}
the second inequality coming from Gelfond's lemma (with an implicit constant depending only on $m$). Choose $j\in\{1,\dots,k\}$ such that $|B_j(\xi)|$ is minimal and fix $i\in\{1,\dots,k\}$ with $i\neq j$. Since $B_i$ and $B_j$ are coprime, their resultant $\Res(B_i,B_j)$ is a non-zero integer. Using Lemma~\ref{lem: estimation classique du resultant}, we find
\begin{align*}
1 \leq |\Res(B_i,B_j)| & \ll \normH{B_i}^{d_j-1} \normH{B_j}^{d_i-1}\big( \normH{B_j}|B_i(\xi)| + \normH{B_i}|B_j(\xi)|\big) \\
& \ll \normH{B_i}^{d_j} \normH{B_j}^{d_i}|B_i(\xi)|,
\end{align*}
with an implicit constant depending only on $\xi$ and $m$, hence
\[
-\log |B_i(\xi)| \leq d_j\log \normH{B_i} + d_i\log \normH{B_j} + \GrO(1).
\]
On the other hand, by summing the above inequalities for $i\neq j$, and by using \eqref{eq proof: log H(P) = sum log H(R_i)}, we obtain
\begin{align*}
\sum_{\substack{i=1 \\ i\neq j}}^{k} -\log |B_i(\xi)|
& \leq d_j\sum_{\substack{i=1 \\ i\neq j}}^{k} \log \normH{B_i} +(m-d_j) \log \normH{B_j} + \GrO(1) \\
& \leq (m-1)\log \normH{R} + \GrO(1).
\end{align*}
We easily deduce that
\[
\prod_{\substack{i=1 \\ i\neq j}}^{k} |B_i(\xi)| \gg \normH{R}^{-(m-1)} \AND
|R(\xi)| = \prod_{i=1}^{k}|B_i(\xi)| \gg |B_j(\xi)| \normH{R}^{- (m-1)}.\qedhere
\]
\end{proof}

\section{A sequence of irreducible polynomials}
\label{section: sequence des Q_i}

As explained in the introduction, to get the upper bound $\homega_n(\xi) \leq 2n-1$, the strategy of Davenport and Schmidt \cite{davenport1969approximation} consists in considering the resultant $\Res(P,Q)$ of two ``good'' polynomial approximations $P,Q\in \bZ[X]_{\leq n}$. To ensure that $\Res(P,Q)$ does not vanish, they need a polynomial $P$ which is irreducible (for it is then easy to find $Q$ so that $P$ and $Q$ are coprime). The same difficulty appears in \cite{bugeaudSchleischitz2016Uniform}. Similarly, we will not work directly with a sequence of minimal polynomials. Instead, we will considerer the largest irreducible factors of the minimal polynomials. Now, let $n,d$ be integers with
\[
2\leq d < 1+ \frac{n}{2}.
\]
In this section, we assume that the transcendental real number $\xi$ satisfies $\homega_n(\xi) > 2n -d$ and we fix a real number $\homega$ (arbitrarily close to $\homega_n(\xi)$) such that
\begin{equation}
\label{eq: notation homega}
\homega_n(\xi) > \homega > 2n -d.
\end{equation}

We denote by $(P_i)_{i\geq 0}$ a sequence of minimal polynomials associated to $n$ and $\xi$. Our goal is to prove the existence of a sequence $(Q_i)_{i\geq 0}$ as below.

\begin{Prop}
\label{prop: existence Qi}
Suppose that \eqref{eq: notation homega} holds. Then, there exist a sequence $(Q_i)_{i\geq 0}$ of pairwise distinct polynomials in $\bZ[X]_{\leq n}$ and an index $j_0\geq 0$ with the following properties. The sequence $(\normH{Q_i})_{i\geq 0}$ is bounded below by $2$, unbounded and non-decreasing, and for any $i\geq 0$
\begin{enumerate}
\item \label{enum: property (Q_i) item 1} $Q_i$ is irreducible (over $\bZ$) and has degree at least $n-d+2$;
\item $Q_i$ divides $P_j$ for some index $j\geq j_0$ (not necessarily unique), and for each $j\geq j_0$ there exists $k\geq 0$ such that $Q_k$ divides $P_j$;
\item \label{enum: property (Q_i) item 6} $|Q_i(\xi)| = \normH{Q_i}^{-\omega(Q_i)}\leq \normH{Q_i}^{-\homega}$, and we further have
\begin{equation}
\label{eq: omega_n donné par les Q_i}
\omega_n(\xi) = \limsup_{k\to \infty} \omega(Q_k) \AND \liminf_{k\to \infty} \omega(Q_k) \geq \homega_n(\xi).
\end{equation}
\item \label{enum: H(P_j) controle par H(Q_i)} if $Q_i$ divides a minimal polynomial $P_j$ with $j\geq j_0$, then
\begin{equation}
\label{eq: premiere estimation Pj vs Qi}
\normH{P_j} \leq \normH{Q_i}^{1+\theta_i}, \quad \textrm{where } \theta_i = \frac{\omega(Q_i)-2n+d}{n-2d+3};
\end{equation}
\item \label{enum: H(Q_i+1) controle par H(Q_i)} we have
\begin{equation}
\label{eq: estimation Q_i+1 par Q_i}
\normH{Q_{i+1}}^\tau \leq \normH{Q_i} \quad \textrm{where } \tau = \frac{\homega\big(\homega -n-d+3\big)}{\omega_n(\xi)\big(\omega_n(\xi)-n-d+3\big)},
\end{equation}
with the convention $\tau= 0$ if $\omega_n(\xi) = \infty$.
\end{enumerate}
\end{Prop}

The above proposition is essentially a consequence of Lemma~\ref{lem: existence facteur irreductible grand} below. Assertion~\eqref{enum: property (Q_i) item 6} ensures that the polynomials $Q_i$ are quite good approximations, and they can be used to compute the exponent of best approximation $\omega_n(\xi)$. Estimate \eqref{eq: estimation Q_i+1 par Q_i} is the analog of the second inequality of \eqref{eq: H(P_i+1) controle par H(P_i)} but is way more difficult to prove. The main reason behind this difficulty is that there may be many polynomials $P\in\bZ[X]_{\leq n}$ with $\normH{Q_i} < \normH{P} < \normH{Q_{i+1}}$ and $|P(\xi)| < Q_i(\xi)$.

In order to prove Proposition~\ref{prop: existence Qi}, we need the two technical lemmas below. Essentially, they will be used to prove that the factors of $P_i$ of small degree are bad approximations. This will lead to the existence of a factor of large degree which is necessarily a rather good approximation.

\begin{Lem}
\label{lemme: omega(R) petit degre}
Suppose that \eqref{eq: notation homega} holds. Then, there exists a constant $c\in(0,1)$ depending only on $\xi$ and $n$ such that for any non-zero polynomial $R\in \bZ[X]_{\leq n-d+1}$ we have
\begin{equation}
\label{lemme; eq1 : omega(R) petit degree}
|R(\xi)| \geq c\normH{R}^{-(n + \deg(R) - 1)} \geq c\normH{R}^{-(2n-d)}.
\end{equation}
In particular \eqref{lemme; eq1 : omega(R) petit degree} holds for any $R\in\bZ[X]_{\leq d-2}$.
\end{Lem}

\begin{proof}
If $R$ is constant we have $|R(\xi)| = \normH{R}$ and the result is trivial. Now, suppose that $R$ is irreducible and not constant. We adapt the arguments of Davenport and Schmidt \cite[\S 5--6]{davenport1969approximation}. Set $H = e^{-n}\normH{R}$. By definition of $\homega_n(\xi)$ and $\homega$, if $H$ is sufficiently large, there exists a non-zero $P\in\bZ[X]_{\leq n}$ such that
\[
\normH{P} \leq H \AND |P(\xi)| \leq H^{-\homega}.
\]
By \eqref{eq: Gelfond's lemma}, the (irreducible) polynomial $R$ is not a factor of $P$, they are thus coprime polynomials. Their resultant is a non-zero integer, and using Lemma~\ref{lem: estimation classique du resultant}, we obtain
\begin{align*}
1 & \ll \normH{P}^{\deg(R)-1}\normH{R}^{n}|P(\xi)| + \normH{P}^{\deg(R)}\normH{R}^{n-1}|R(\xi)| \\
& \ll H^{n+\deg(R)-1-\homega} + H^{n+\deg(R)-1}|R(\xi)|.
\end{align*}
Since $\homega > 2n-d$ and $\deg(R) \leq n-d+1$, the first term tends to $0$ as $H$ tends to infinity. Hence $ 1 \ll H^{n+\deg(R)-1}|R(\xi)|$, which implies \eqref{lemme; eq1 : omega(R) petit degree}.

If $R$ is not irreducible, we write $R = \prod_{i=1}^{s}R_i$ with integer $s\geq 1$ and $R_1,\dots,R_s\in\bZ[X]$ irreducible of degree $\leq \deg(R)$ (possibly constant). Combining $\normH{R} \asymp \prod_{i=1}^{s}\normH{R_i}$ together with \eqref{lemme; eq1 : omega(R) petit degree} applied to the irreducible factors $R_i$, we find
\begin{equation*}
|R(\xi)| = \prod_{i=1}^{s}|R_i(\xi)| \gg \prod_{i=1}^{s}\normH{R_i}^{-(n + \deg(R) - 1)} \gg \normH{R}^{-(n + \deg(R) - 1)}.
\end{equation*}
Finally, the last assertion comes from the fact that $d-1\leq n+d-1$ (since $d\leq 1+n/2$).
\end{proof}

\begin{Lem}
\label{lem: existence facteur irreductible grand}
Suppose that \eqref{eq: notation homega} holds. There exist $i_0\geq 0$ and a constant $c>0$ such that for each $i\geq i_0$ the polynomial $P_i$ has a unique irreducible factor $\tP_i\in\bZ[X]$ of degree $\geq n-d+2$ and positive leading coefficient. It satisfies
\begin{equation}
\label{eq lem: existence facteur irr grand eq 2}
|P_i(\xi)| \normH{P_i}^{n+d-3} \geq c|\tP_i(\xi)| \normH{\tP_i}^{n+d-3},
\end{equation}
moreover $\big(\normH{\tP_i}\big)_{i\geq i_0}$ tends to infinity and as $i$ tends to infinity. For each $i$ large enough, we have $\normH{\tP_i}>1$, and writing $|\tP_i(\xi)| = \normH{\tP_i}^{-\omega(\tP_i)}$, we furthermore have
\begin{equation}
\label{eq lem: existence facteur irr grand eq 1}
\omega_n(\xi) = \limsup_{i\to\infty} \omega(\tP_i) \AND \liminf_{i\to\infty} \omega(\tP_i) \geq \homega_n(\xi).
\end{equation}
\end{Lem}

\begin{proof}
First, note that since $d < 1+n/2$, if we decompose $P_i$ as a product of irreducibles, there is at most one factor of degree $\geq n-d+2$. Fix $i\geq 0$ large enough so that $\omega(P_i) \geq \homega$, and write
\[
P:=P_i = \prod_{k=1}^s R_k,
\]
where $R_1,\dots,R_s\in\bZ[X]$ are irreducible polynomials (and $s$ is a positive integer). Suppose that $\deg(R_k) \leq n-d+1$ for each $k=1,\dots,s$. Then, by Lemma~\ref{lemme: omega(R) petit degre} together with $\normH{P} \asymp \prod_k \normH{R_k}$, we find
\begin{equation*}
\normH{P}^{-\homega}\geq |P(\xi)| = \prod_{k=1}^{s} |R_k(\xi)| \gg \prod_{k=1}^{s} \normH{R_k}^{-(2n-d)} \asymp \normH{P}^{-(2n-d)}.
\end{equation*}
This is impossible if $i$ is sufficiently large since $\homega > 2n-d$. Therefore, if $i$ is large enough, one of the factors $R_k$ has degree at least $n-d+2$. Without loss of generality, we may suppose that it is $R:=R_1$.
Write $S:= \prod_{k = 2}^{s} R_k$, so that $P=RS$. We have $\deg(S) \leq d-2$, and \eqref{lemme; eq1 : omega(R) petit degree} of Lemma~\ref{lemme: omega(R) petit degre} yields
\begin{equation*}
|S(\xi)| \gg \normH{S}^{-(n+d-3)}.
\end{equation*}
Together with $\normH{P}\asymp \normH{R}\normH{S}$, this leads to
\begin{equation*}
|P(\xi)| = |R(\xi)||S(\xi)| \gg |R(\xi)|\normH{S}^{-(n+d-3)} \asymp |R(\xi)|\normH{R}^{n+d-3}\normH{P}^{-(n+d-3)},
\end{equation*}
and \eqref{eq lem: existence facteur irr grand eq 2} follows easily by setting $\tP_i:=R$. The rest of the proof is based only on \eqref{eq lem: existence facteur irr grand eq 2} and the inequality $\normH{\tP_i} \ll \normH{P_i}$. Note that $|P_i(\xi)|\normH{P_i}^{n+d-3} \ll \normH{P_i}^{n+d-3-\homega}$ tends to $0$ as $i$ tends to infinity (using $d < 1+n/2$ together with $\homega > 2n-d$). We deduce that $|\tP_i(\xi)|\normH{\tP_i}^{n+d-3}$ also tends to $0$ as $i$ tends to infinity, which is possible only if $\normH{\tP_i}$ tends to infinity. In particular, if $i$ is large enough we must have $\normH{\tP_i} > 1$. Writing $|\tP_i(\xi)| = \normH{\tP_i}^{-\omega(\tP_i)}$, we also have $\omega(\tP_i) > n+d-3$. Now, using $\normH{\tP_i}\ll \normH{P_i}$, and taking the logarithms of the two sides of \eqref{eq lem: existence facteur irr grand eq 2}, we get
\begin{align*}
\big(\omega(P_i)-(n+d-3)\big)\log \normH{P_i} & \leq \big(\omega(\tP_i)-(n+d-3)\big)\log \normH{\tP_i} + \GrO(1) \\
& \leq \big(\omega(\tP_i)-(n+d-3)\big)\big(\log \normH{P_i}+\GrO(1)\big) + \GrO(1).
\end{align*}
By dividing by $\log \normH{P_i}$ and by simplifying, we deduce that $\omega(\tP_i) \geq \omega(P_i)\big (1 - o(1)\big)$ and \eqref{eq lem: existence facteur irr grand eq 1} follows easily from \eqref{eq: encadrement omega(P_i)}.
\end{proof}

\begin{proof}[Proof of Proposition~\ref{prop: existence Qi}]
Let $i_0\geq 0$ and $(\tP_i)_{i\geq i_0}$ given by Lemma~\ref{lem: existence facteur irreductible grand}. Let $(Q_i)_{i\geq 0}$ be the (infinite) sequence of factors $(\tP_j)_{j\geq i_0}$ reordered by increasing height, without repetition. By Lemma~\ref{lem: existence facteur irreductible grand}, we may assume $i_0$ large enough so that $\normH{Q_i} > 1$ for each~$i$, as well as $|Q_i(\xi)| \leq \normH{Q_i}^{-\homega}$. This sequence clearly satisfies the first assertions \eqref{enum: property (Q_i) item 1} to \eqref{enum: property (Q_i) item 6}, the third one coming from \eqref{eq lem: existence facteur irr grand eq 1} together with \eqref{eq: encadrement omega(P_i)}.

Now, let $i\geq 0$ and let $j\geq i_0$ be an index such that $Q_i$ divides $P_j$. Since we have $\normH{Q_i}\ll \normH{P_j}$ by Gelfond's lemma, the index $j$ tends to infinity as $i$ tends to infinity. Then, estimate \eqref{eq lem: existence facteur irr grand eq 2} can be rewritten as
\begin{equation}
\label{eq proof: inter 1 pour controler H(P_j)}
|P_j(\xi)|^{-1}\normH{P_j}^{-n-d+3} \ll |Q_i(\xi)|^{-1}\normH{Q_i}^{-n-d+3} = \normH{Q_i}^{\omega(Q_i)-n-d+3}.
\end{equation}
Using $|P_j(\xi)|^{-1} \gg \normH{P_j}^{\homega}$ and $\homega > 2n-d$, we get, for each large enough $i$,
\begin{equation*}
\normH{P_j}^{n-2d+3} \leq \normH{Q_i}^{\omega(Q_i)-n-d+3},
\end{equation*}
which is equivalent to \eqref{eq: premiere estimation Pj vs Qi}. So, assertion~\eqref{enum: H(P_j) controle par H(Q_i)} holds assuming $i_0$ large enough.

It remains to prove assertion~\eqref{enum: H(Q_i+1) controle par H(Q_i)}. Note that this is trivial if $\omega_n(\xi) = \infty$. Let us assume that $\omega_n(\xi) < \infty$ and fix a small $\ee > 0$ to be chosen later. For each pair $(i,j)$ as above with $j \geq i_0$ large enough as a function of $\ee$, we have $\omega(P_j) > \homega_n(\xi)-\ee/2$ and $\omega(Q_i) < \omega_n(\xi)+\ee/2$, and thus \eqref{eq proof: inter 1 pour controler H(P_j)} yields
\begin{equation*}
\normH{P_j}^{\homega_n(\xi)-\ee-n-d+3} \leq \normH{Q_i}^{\omega_n(\xi)+\ee-n-d+3},
\end{equation*}
for each $i \geq 0$ and each $j\geq i_0$ such that $Q_i$ divides $P_j$. We define $k$ as the largest index such that
\begin{equation*}
\normH{P_k} \leq \normH{Q_i}^{\theta(\ee)},\quad \textrm{where } \theta(\ee) = \frac{\omega_n(\xi)+\ee-n-d+3}{\homega_n(\xi)-\ee-n-d+3}.
\end{equation*}
Since $\normH{P_j} \leq \normH{Q_i}^{\theta(\ee)}$, by maximality of $k$ we have $i_0\leq j\leq k$. Let $\ell$ be such that $Q_\ell$ divides $P_{k+1}$. We find
\begin{equation*}
\normH{P_k} \leq \normH{Q_i}^{\theta(\ee)} < \normH{P_{k+1}} \leq \normH{Q_\ell}^{\theta(\ee)},
\end{equation*}
and therefore $\ell \geq i+1$. On the other hand, since by Gelfond's lemma we have $\normH{Q_\ell} \ll \normH{P_{k+1}}$, we deduce from \eqref{eq: H(P_i+1) controle par H(P_i)} that
\begin{equation*}
\normH{Q_{i+1}} \leq \normH{Q_\ell} \ll \normH{P_{k+1}} \ll \normH{P_k}^{\omega_n(\xi)/\homega}\leq \normH{Q_i}^{\omega_n(\xi)\theta(\ee) /\homega}.
\end{equation*}
We now choose $\ee>0$ small enough so that
\begin{equation*}
\theta(\ee) < \frac{\omega_n(\xi)-n-d+3}{\homega -n-d+3}.
\end{equation*}
This is possible since $\homega < \homega_n(\xi)$, and it yields \eqref{eq: estimation Q_i+1 par Q_i} for each $i\geq 0$, assuming that $i_0$ is large enough.
\end{proof}

\section{On the dimension of some polynomial subspaces}
\label{section: espaces V_N}

We start by introducing some families of vector spaces spanned by polynomials, and we study their dimensions.

\begin{Def}
\label{Def: fonction g_A(N)}
Let $k\geq n$ be an integer and let $\cA$ be a subset of $\bR[X]_{\leq n}$. We~define
\begin{align*}
\cB_k(\cA) &= \big\{ Q, XQ,\dots, X^{k-\deg(Q)}Q \,;\, Q\in\cA\setminus\{0\} \big\} \subset \bR[X]_{\leq k}, \\
V_k(\cA) &= \Vect[\bR]{\cB_k(\cA)}, \\
g_\cA(k) &= \dim V_k(\cA).
\end{align*}
\end{Def}

The spaces $V_k(\cA)$ play the role of the spaces $\cU^{k}(\cA)$ in \cite[\S3]{poels2022simultaneous} (for simultaneous approximation). We obtain analog properties. Note that if $\cA$ contains at least one non-zero polynomial, then
\begin{equation}
\label{eq: suite V_k(A) strict. croissante}
V_n(\cA) \varsubsetneq V_{n+1}(\cA) \varsubsetneq\cdots.
\end{equation}

The goal of this section is to prove the following result. We could not find a reference for the proposition below.

\begin{Prop}
\label{Cor: V_2n-k = tout l'espace}
Let $k$ be an integer with $0\leq k \leq n$, and let $\cA$ be a set of $k+1$ linearly polynomials of $\bR[X]_{\leq n}$. Suppose that the $\gcd$ of the elements of $\cA$ is $1$ (in~other words, the ideal spanned by $\cA$ is $\bR[X]$). Then
\begin{equation}
\label{Cor: V_N = tout l'espace}
V_{2n-k}(\cA) = \bR[X]_{\leq 2n-k}.
\end{equation}
\end{Prop}

The case $k=1$ is a classical result (it is implied by the fact that the resultant of two coprime polynomials is non-zero). The proof of Proposition~\ref{Cor: V_2n-k = tout l'espace} is given at the end of the section. Recall that a function $f : \{n,n+1, \ldots \} \to \bR$ is \emph{concave} if for any $i > n$, it satisfies
\[
f(i)-f(i-1) \geq f(i+1)-f(i).
\]
The next result is a dual version of \cite[Prop.\,3.1]{poels2022simultaneous} (where we deal with simultaneous approximation to the successive powers of $\xi$).

\begin{Lem}
Let $\cA \neq \{0\}$ be a non-empty subset of $\bR[X]_{\leq n}$. The function $g_\cA$ is concave and (strictly) increasing on $\{n,n+1,\ldots\}$.
\end{Lem}

\begin{proof}
The series of inclusions \eqref{eq: suite V_k(A) strict. croissante} shows that the function $g_\cA$ is increasing on $\{n,n+1,\dots\}$. For simplicity, we write $V_i=V_{i}(\cA)$ and $\cB_{i} = \cB_{i}(\cA)$ for each $i \geq n$. Given an integer $i\geq n$ we have $XV_{i}\subset V_{i+1}$, and we set
\[
h(i):= \dim \big(V_{i+1}/XV_{i}\big) = g_\cA(i+1)-g_\cA(i).
\]
We have to prove that $h$ is decreasing on $\{n,n+1,\cdots\}$. Fix $i\geq n+1$ and consider the linear map $\pi : V_i \to V_{i+1}/XV_{i}$ defined by $\pi(P) = P + XV_{i}$. Since $\cB_{i}\cup X\cB_{i} = \cB_{i+1}$, we have $V_i + XV_i = V_{i+1}$. So $\pi$ is surjective, and consequently $\textrm{Im } \pi = V_{i+1}/XV_{i}$ is isomorphic to $V_i / \ker \pi$.
On the other hand, $XV_{i-1}\subset V_{i} \cap XV_{i} \subset \ker\pi$, so $XV_{i-1}$ is subspace of $\ker \pi$. Hence
\begin{equation*}
h(i-1) = \dim \big(V_{i}/XV_{i-1}\big) \geq \dim\big(V_i / \ker \pi \big) = \dim \big(V_{i+1}/XV_{i}\big) = h(i).\qedhere
\end{equation*}
\end{proof}

\begin{Lem}
\label{lem : dim V_n+j(P,Q)}
Let $P,Q \in \bR[X]_{\leq n}$ be two coprime polynomials. Then, we have
\begin{equation*}
\dim V_{n+j}(P,Q) \geq 2(j+1),
\end{equation*}
for each $j\in\{0,\dots,n-1\}$. In particular $V_{2n-1}(P,Q) = \bR[X]_{\leq 2n-1}$.
\end{Lem}

\begin{proof}
Let $p$ (\resp $q$) denote the degree of $P$ (\resp of $Q$). There exist $\alpha,\beta\in\bR$ such that the polynomial $\tP := P(X)(X-\alpha)^{n-p}$ and $\tQ:= Q(X)(X-\beta)^{n-q}$ are coprime (and of degree exactly $n$). Fix $j\in\{0,\dots,n-1\}$. Since $\tP$ and $\tQ$ are coprime and $j<n$, the linear map
\begin{align*}
\bR[X]_{\leq j}\times \bR[X]_{\leq j} & \longrightarrow \bR[X]_{\leq n+j} \\
(R,S) & \longmapsto R\tP + S\tQ
\end{align*}
is injective, so its image $V_{n+j}(\tP,\tQ) \subset V_{n+j}(P,Q)$ has dimension $2(j+1)$.
\end{proof}

\begin{proof}[Proof of Proposition~\ref{Cor: V_2n-k = tout l'espace}]
For simplicity, we write $g=g_\cA$. Recall that $\cA$ has cardinality $k+1$, so that $g(n)\geq \textrm{card}(\cA) = k+1$. If $k=n$, then \eqref{Cor: V_N = tout l'espace} is automatic (since in that case $\cA$ contains a basis of $\bR[X]_{\leq n}$). So, we may assume that $k< n$. We first prove that for each sufficiently large $m$, we have
\begin{equation}
\label{eq : il existe m tel que V_m = tout}
V_{m}(\cA) = \bR[X]_{\leq m}.
\end{equation}
Indeed, since the ideal spanned by $\cA$ is $\bR[X]$, there exists an integer $\ell\geq n$ such that $1\in V_\ell(\cA)$. Let $P$ be a non-zero element in $\cA$ of degree $d$, and set $m=\ell+d$. Then $V_{m}(\cA)$ contains $\bR[X]_{\leq d}$, as well as the polynomials $P,XP,\ldots,X^\ell P$. We easily deduce \eqref{eq : il existe m tel que V_m = tout}.

By contradiction, suppose that \eqref{Cor: V_N = tout l'espace} does not hold, \ie
\begin{equation}
\label{eq proof:Cor: V_N = tout l'espace: by contradiction}
g(2n-k) \leq 2n-k.
\end{equation}
We distinguish two cases. Suppose first that $g(2n-k)-g(2n-k-1) \geq 2$. By concavity, then $g(j)-g(j-1)\geq 2$ for each $j$ with $n < j \leq 2n-k$, and we deduce that
\begin{equation*}
g(2n-k) \geq g(n)+2(n-k) \geq k+1 + 2(n-k) = 2n-k+1,
\end{equation*}
since $g(n) \geq \textrm{card}(\cA) = k+1$. This contradicts \eqref{eq proof:Cor: V_N = tout l'espace: by contradiction}, so $g(2n-k)-g(2n-k-1) \leq 1$. Since the function $g$ is increasing and concave, it is linear with slope $1$ on
\[
\{2n-k,2n-k+1,\dots\}.
\]
Choosing $m > 2n-k$ such that \eqref{eq : il existe m tel que V_m = tout} holds, we obtain by \eqref{eq proof:Cor: V_N = tout l'espace: by contradiction}
\begin{equation*}
m+1 = g(m) = g(2n-k)+ m - (2n-k) \leq m,
\end{equation*}
a contradiction. Hence \eqref{Cor: V_N = tout l'espace} holds.
\end{proof}

\begin{figure}[htb]
\begin{tikzpicture}[xscale=0.5,yscale=0.4, line cap=round,line join=round,>=triangle 45,x=1cm,y=1cm]
\clip(-5,-2) rectangle (20,15);
\draw[-stealth, semithick] (-0.15,0)--(18,0) node[below]{$i$};
\draw[-stealth, semithick] (0,-0.15)--(0,14.5) node[left]{$g(i)$};

\draw [thick] (0,1)-- (1,5);
\draw [thick] (1,5)-- (2.66,8.059);
\draw [thick] (2.66,8.059)-- (7,10.5);
\draw [thick] (7,10.5)-- (14,12);
\draw [dashed] (7,13.34)-- (7,-0.5);

\node[left] at (-0,1.7) {$\geq \textrm{card}(A)$};
\node[above] at (4,12) {slope $\geq 2$};
\node[above] at (10,12) {slope $= 1$};
\node[below] at (7,0) {$2n-\ell$};
\node[below] at (14,0) {$2n-1$};
\node[left] at (0,13) {$2n$};

\draw [dash pattern=on 2pt off 3pt](0,12)-- (14,12);
\draw [dash pattern=on 2pt off 3pt] (14,12)-- (14,0);

\node[below] at (0,-0.15) {$i=n$};
\node[draw,circle,inner sep=1.25pt,fill] at (0,1) {};
\node[draw,circle,inner sep=1.25pt,fill] at (1,5) {};
\node[draw,circle,inner sep=1.25pt,fill] at (2.666,8.059) {};
\node[draw,circle,inner sep=1.25pt,fill] at (7,10.5) {};
\node[draw,circle,inner sep=1.25pt,fill] at (14,12) {};
\node[draw,circle,inner sep=1.25pt,fill] at (0,12) {};
\node[draw,circle,inner sep=1.25pt,fill] at (14,0) {};
\node[draw,circle,inner sep=1.25pt,fill] at (7,0) {};
\node[draw,circle,inner sep=1.25pt,fill] at (0,0) {};
\end{tikzpicture}
\caption{Graph of the piecewise linear function interpolating the
values $g(i)=\dim V_i(\cA)$ at integers $i\in\{n,\dots,2n-1\}$.}
\label{fig:1}
\end{figure}

\section{Proof of Theorem~\ref{Thm : main d=2} (case \texorpdfstring{$d=2$}{d2})}
\label{Section: cas d = 2}

In this section, we deal with the case $d=2$ to prove Theorem~\ref{Thm : main d=2}, namely that $\homega_3(\xi)\leq 2+\sqrt 5 = 4.23\cdots$ and $\homega_n(\xi) \leq 2n-2$ for each $n\geq 4$. The estimate $\homega_n(\xi) \leq 2n-2$ was already known for $n\geq 10$, however for $n=4,\dots,9$ it is a new result. For $n=3$, our bound improves on the bound $\homega_3(\xi)\leq 3+\sqrt 2 = 4.41\cdots$ due to Bugeaud and Schleischitz \cite{bugeaudSchleischitz2016Uniform}. Moreover, our proof does not require Marnat-Moshchevitin's inequality \cite{marnat2018optimal}.

\begin{proof}[Proof of Theorem~\ref{Thm : main d=2}]
Suppose that $\homega_n(\xi) > 2n-2$, and fix a real number $\homega$ such that
\[
\homega_n(\xi) > \homega > 2n-2.
\]
Let $(P_i)_{i\geq 0}$ be a sequence of minimal polynomials associated to $n$ and $\xi$ as in Section~\ref{section: minimal pol}. According to Lemma~\ref{lem: existence facteur irreductible grand} (with $d=2$) there exists an index $i_0\geq 0$ such that $P_i$ has degree $n$ and is irreducible for each $i\geq i_0$. Consequently, up to a finite number of terms, the sequence $(P_i)_{i\geq 0}$ coincides with the sequence $(Q_i)_{i\geq 0}$ of Proposition~\ref{prop: existence Qi}. Let $I$ denotes the set of indices $i\geq i_0+1$ such that $P_{i-1}$, $P_i$ and $P_{i+1}$ are linearly independent. By Lemmas~\ref{Lem: P_iH_i+1 = P_j-1H_j} and~\ref{lem: I infini}, the set $I$ is infinite, and for any consecutive $i<j$ in $I$, we have
\[
\normH{P_{i+1}}|P_i(\xi)| \asymp \normH{P_j}|P_{j-1}(\xi)|.
\]
Furthermore, the irreducible polynomials $P_i$ and $P_{i+1}$ are also coprime since $\normH{P_i} < \normH{P_{i+1}}$ and $|P_i(\xi)| > |P_{i+1}(\xi)|$. Lemma~\ref{lem: estimation classique du resultant} yields
\begin{align*}
1 \leq |\Res(P_i,P_{i+1})| \ll \normH{P_i}^{n-1}\normH{P_{i+1}}^n|P_{i}(\xi)| &\ll \normH{P_i}^{n-1}\normH{P_{j}}^n|P_{j-1}(\xi)| \\
& \ll \normH{P_i}^{n-1}\normH{P_{j}}^{n-\homega}.
\end{align*}
We deduce that
\begin{equation}
\label{eq proof: estimation theta}
\normH{P_{j}} \leq \normH{P_i}^\theta \quad \textrm{where } \theta = \frac{n-1}{\homega-n}.
\end{equation}
Let $h<i<j$ be consecutive indices in $I$. We have the following configuration
\[
\Vect[\bR]{P_h,P_{h+1}} = \Vect[\bR]{P_{i-1},P_{i}} \neq \Vect[\bR]{P_{i},P_{i+1}},
\]
so $P_{h}, P_{h+1}, P_{i+1}$ are linearly independent. Proposition~\ref{Cor: V_2n-k = tout l'espace} combined with Lemma~\ref{lem : dim V_n+j(P,Q)} implies that
\begin{equation*}
\Big(\bR[X]_{\leq n-2}P_{h}\oplus \bR[X]_{\leq n-2}P_{h+1}\Big) + \bR[X]_{\leq n-2}P_{i+1} = \bR[X]_{\leq 2n-2}.
\end{equation*}
Choose $k\in\{0,\dots,n-2\}$ such that $\big(P_{h},\dots,X^{n-2}P_{h},P_{h+1},\dots,X^{n-2}P_{h+1}, X^kP_{i+1}\big)$ is a basis of $\bR[X]_{\leq 2n-2}$. We denote by $M$ the matrix of this basis expressed in the canonical basis $(1,X,\dots,X^{2n-2})$. Estimating $\det(M)$ as in the proof of Lemma~\ref{lem: estimation classique du resultant} (in other words, for $\ell=2,\dots,2n-2$, we add to the first row of $M$ the $\ell$-th row multiplied by $\xi^{\ell-1}$), we get the estimates
\begin{equation*}
1 \leq \big|\det\big(M)\big| \ll |P_{h}(\xi)|\normH{P_{h}}^{n-2}\normH{P_{h+1}}^{n-1}\normH{P_{i+1}}.
\end{equation*}
Now, since $\normH{P_{h+1}}^{n-1}|P_{h}(\xi)| \asymp \normH{P_{h+1}}^{n-2}\normH{P_i}|P_{i-1}(\xi)| \ll \normH{P_i}^{n-1-\homega}$, we deduce that
\begin{equation}
\label{eq proof: estimation tau}
\normH{P_i}^{\homega-n+1} \ll \normH{P_{h}}^{n-2}\normH{P_{j}}.
\end{equation}
For consecutive $i<j$ in $I$, define $\tau_i\in(0,1)$ by
\[
\normH{P_i} = \normH{P_j}^{\tau_i}
\]
and set $\tau = \limsup_{i\in I, i\to\infty} \tau_i \in [0,1]$. Let $h<i<j$ be consecutive indices in $I$ as previously. By \eqref{eq proof: estimation tau}, we obtain
\[
\homega-n+1 \leq (n-2)\tau_h + \frac{1}{\tau_i} +o(1) \leq (n-2)\tau + \frac{1}{\tau_i} +o(1).
\]
We infer that
\begin{equation}
\label{eq proof: majoration homega_3: eq 2}
p(\tau) \geq 0, \quad \textrm{where } p(t) = (n-2)t^2-(\homega-n+1)t+1.
\end{equation}
Note that
\begin{equation*}
p(0) = 1, \qquad p\Bigl(\frac{1}{n-2}\Bigr) = \frac{2n-2-\homega}{n-2} < 0 \AND p(1) = 2n-2-\homega < 0.
\end{equation*}
We deduce that $p$ has one root $\alpha\in(0,1/(n-2))$ and one root larger than $1$. Since $\tau\in[0,1]$ and $p(\tau)\geq 0$, we obtain $\tau \leq \alpha$. Combined with the estimate $\normH{P_i} = \normH{P_j}^{\tau_i} \ll \normH{P_i}^{\theta\tau_i}$ valid
for any $i\in I$ (this is a consequence of \eqref{eq proof: estimation theta}), this leads to
\begin{equation}
\label{eq proof: majoration homega_3: eq 3}
1 \leq \theta\tau \leq \theta\alpha < \frac{n-1}{(n-2)^2}.
\end{equation}
We easily check that this is impossible when $n\geq 4$ (the right-hand side is strictly less than $1$), thus $\homega_n(\xi)\leq 2n-2$ for each $n\geq 4$.

We now deal with the case $n=3$. Suppose by contradiction that $\homega_3(\xi) > 2+\sqrt 5$ and choose $\homega$ such that
\[
\homega_3(\xi) > \homega > 2+\sqrt 5.
\]
The polynomial $p$ from \eqref{eq proof: majoration homega_3: eq 2} becomes $p(t) = t^2-(\homega-2)t+1$. Denote by $\alpha$ its smallest root, and by $\beta = (\sqrt 5 -1)/2$ the smallest root of the polynomial $t^2-\sqrt 5 t + 1$. We find
\begin{equation*}
0 = \beta^2-\sqrt 5 \beta + 1 > \beta^2-(\homega-2 )\beta + 1 = p(\beta),
\end{equation*}
hence $\alpha < \beta$. Combined with $\theta = 2/(\homega-3) < 1/\beta$, this implies that $\theta \alpha < 1$, which contradicts \eqref{eq proof: majoration homega_3: eq 3}. It follows that $\homega_3(\xi) \leq 2+\sqrt 5$.
\end{proof}

\section{Multilinear algebra and height of polynomial subspaces}

This section is divided into two parts. We introduce and study a quantity $\HHstar(V)$ associated to a subspace $V\subset\bR^m$ defined over $\bQ$ in Section~\ref{def petitesse HHstar(V)}. Intuitively, $\HHstar(V)$ is small if $V$ is spanned by good polynomials approximations of $\bZ[X]$ (\ie small when evaluated at $\xi$). This will be a key-point for estimating the height of the polynomials~$Q_i$ of Section~\ref{section: sequence des Q_i}. In order to define $\HHstar$, we need some tools of multilinear algebra that we recall in Section~\ref{section: hauteurs tordues}. In the appendix we give another interpretation of $\HHstar$ in term of twisted heights.

\subsection{Multilinear algebra and Hodge duality}
\label{section: hauteurs tordues}

For each integer $m$, we view $\bR^{m+1}$ as an Euclidean space for the usual scalar product $\psc{\cdot}{\cdot}$, and we denote by $\norm{\cdot}$ the associated Euclidean norm. For each $k = 1, \cdots , m+1$, we identify $\sbw{k}{\bR^{m+1}}$ with $\bR^N$, where $N = \binom{m+1}{k}$, via an ordering of the Plücker coordinates, and we denote by $\norm{\by}$ the norm of a point $\by \in \sbw{k}{\bR^{m+1}}\cong \bR^N$. This is independent of the ordering. Let $V$ be a $k$-dimensional subspace of $\bR^{m+1}$ defined over $\bQ$, \ie such that $\Vect[\bR]{V\cap\bQ^{m+1}} = V$. Its (standard) \emph{height} $\HH(V)$ is the covolume of the lattice $V\cap\bZ^{m+1}$ inside $V$ (with the convention that $\HH(V) = 1$ if $V=\{0\}$). Explicitly, we have
\[
\HH(V) := \norm{ \bx_1\wedge \cdots \wedge \bx_k},
\]
for any $\bZ$-basis $(\bx_1,\dots,\bx_k)$ of the lattice $V\cap \bZ^{m+1}$. Schmidt established the very nice inequality
\begin{equation*}
\HH(U\cap V)\HH(U+V) \leq \HH(U)\HH(V),
\end{equation*}
valid for any subspaces $U,V$ of $\bR^{m+1}$ defined over $\bQ$ (see \cite[Ch.\,I, Lem.\,8A]{Schmidt}). In~this paper, we need to work with a ``twisted'' height and the corresponding version of Schmidt's inequality (obtained by following Schmidt's original arguments).

Let $(\be_1,\dots,\be_{m+1})$ denotes the canonical basis of $\bR^{m+1}$, and let $k$ be an integer with $0\leq k \leq m+1$. The Hodge star operator
\begin{equation*}
* : \sbw{k}{\bR^{m+1}} \mathop{\longrightarrow}^{\sim} \sbw{m+1-k}{\bR^{m+1}}
\end{equation*}
is defined by
\begin{equation*}
*(\be_{i_1}\wedge \cdots \wedge \be_{i_k}) = \ee_{i_1,\dots,i_k}\be_{j_1}\,\wedge \cdots \wedge \be_{j_{m+1-k}}
\end{equation*}
for any indices $i_1 < \cdots < i_k$ and $j_1 < \cdots < j_{m+1-k}$ forming a partition of $\{1,\dots,m+1\}$, where $\ee_{i_1,\dots,i_k}$ denotes the signature of the substitution $(1,\dots,m+1) \mapsto (j_1,\dots,j_{m+1-k}, i_1,\dots,i_k)$. Given $\bX\in \sbw{k}{\bR^{m+1}}$, the point $*\bX$ is called the \emph{Hodge dual} of $\bX$.

We now collect some useful properties of the Hodge star operator, see for example \cite{Hodge}, \cite{bourbakiAlgebre} and \cite[\S 3]{bugeaud2010transfer} for more details. First,
\[
\norm{*\bX} = \norm{\bX} \AND *(*\bX) = (-1)^{k(m+1-k)}\bX
\]
for any $\bX\in \sbw{k}{\bR^{m+1}}$. If $\bX = \bx_1\wedge \cdots \wedge \bx_k$ is a system of Plücker coordinates of a $k$-dimensional subspace $V\subset \bR^{m+1}$, then $*\bX$ is a system of Plücker coordinates of its orthogonal $V^\perp$. This implies the classical identity
\[
\HH(V) = \HH(V^\perp).
\]
If $k>0$, then given $\by\in\bR^{m+1}$ and a multivector $\bX\in \sbw{k}{\bR^{m+1}}$, the point
\begin{equation*}
\by \lrcorner\, \bX = *\big(\by\wedge (*\bX)\big) \in \sbw{k-1}{\bR^{m+1}}
\end{equation*}
is called the \emph{contraction} of $\bX$ by $\by$ (see \cite[Lem.\,2]{bugeaud2010transfer}). Explicitly, if $\bX = \bx_1\wedge \cdots \wedge \bx_k$ is a decomposable multivector, then
\begin{equation}
\label{eq: contraction formule explicite}
\by \lrcorner\, \bX = \sum_{i=1}^{k} (-1)^{k-i} \psc{\bx_i}{\by} \, \bx_1\wedge \cdots \wedge \widehat{\bx_i} \wedge \cdots \wedge \bx_k,
\end{equation}
where the hat on $\bx_i$ means that this term is omitted from the wedge product (see \cite[Eq.\,(3.3)]{bugeaud2010transfer}). In particular, if $k=1$ and $\bX = \bx\in\bR^{m+1}$, we simply have
\begin{equation}
\label{eq: contraction et produit scalaire}
\by \lrcorner\, \bx = \psc{\by}{\bx}.
\end{equation}

\subsection{Schmidt's inequality}
\label{def petitesse HHstar(V)}

Let $m$ be a non-negative integer and set $\Xi_m = (1,\xi,\xi^2,\dots,\xi^m)$. We keep the notation of Section~\ref{section: hauteurs tordues}.

\begin{Def}
\label{Def: def HHstar}
Let $V$ be a $k$-dimensional subspace of $\bR^{m+1}$ defined over $\bQ$, with $k\geq 1$, and let $(\bx_1,\dots,\bx_k)$ be a $\bZ$-basis of the lattice $V\cap \bZ^{m+1}$. We set
\begin{equation*}
\HHstar(V) = \norm{\Xi_m \, \lrcorner\, \bX } = \norm{\Xi_m \wedge (*\bX)},
\end{equation*}
where $\bX = \bx_1\wedge \cdots \wedge\bx_k$. By convention, we set $\HHstar(\{0\}) = 0$. Following the notation of \cite[\S11]{poels2022simultaneous}, we also set
\begin{equation*}
L_\xi(V) = \norm{\Xi_m \wedge \bX},
\end{equation*}
with the convention that $L_\xi(\{0\}) = \norm{\Xi_m}$.
\end{Def}

\begin{remark*}
If $(\bx_1',\dots,\bx_k')$ is another $\bZ$-basis of $V\cap \bZ^{m+1}$, then $\bx_1'\wedge \cdots \wedge\bx_k' = \pm \bX$. Consequently, $\HHstar(V)$ and $L_\xi(V)$ do not depend on the choice of the basis. In \cite{poels2022simultaneous}, we considered $L_\xi(V)$ for spaces $V$ spanned by good simultaneous approximations. The function $\HHstar$ is connected to the quantity introduced in \cite[Def.\,7.1]{poelsroy2022parametric} (where we work in a number field $K$ instead of $\bQ$). Note that $\HHstar(V) = 0$ if and only if $\Xi_m\in V^\perp$. Since $\xi$ is transcendental, this is only possible when $V=\{0\}$. We have
\begin{equation*}
\HHstar(\bR^{m+1}) = \norm{\Xi_m } \asymp 1,
\end{equation*}
where the implicit constants depend on $\xi$ and $m$ only. Moreover, \eqref{eq: contraction et produit scalaire} implies that
\begin{equation}
\label{eq: HHstar(bx)}
\HHstar\big(\Vect[\bR]{\bx}\big) = |\psc{\Xi_m}{\bx}|
\end{equation}
for any primitive integer point $\bx\in\bZ^{m+1}$. Equation \eqref{eq: contraction formule explicite} yields the explicit formula
\begin{equation}
\label{eq: HHstar(V) formule explicite}
\HHstar(V) = \bnorm{ \sum_{i=1}^{k} (-1)^{k-i}\psc{\bx_i}{\Xi_m} \, \bx_1\wedge \cdots \wedge \widehat{\bx_i} \wedge \cdots \wedge \bx_k }.
\end{equation}
On the other hand, if $(\by_1,\dots,\by_{m+1-k})$ is a $\bZ$-basis of $V^\perp\cap \bZ^{m+1}$, then $*\bX = \pm \by_1\wedge \cdots \wedge \by_{m+1-k}$. Consequently, we can also write
\begin{equation}
\label{eq: D_xi(V) = L_xi(V^perp)}
\HHstar(V) = \norm{\Xi_m \wedge \by_1\wedge\cdots \wedge \by_{m+1-k}} = L_\xi(V^\perp).
\end{equation}
Both formulas for $\HHstar(V) $ will be useful.
\end{remark*}

\begin{Prop}[Schmidt's inequality]
\label{Prop: Schmidt inegalite pour HHstar}
For any subspaces $U,V$ of $\bR^{m+1}$ defined over $\bQ$, we have
\begin{equation}
\label{eq: Schmidt's inequality HHstar}
\HHstar(U\cap V)\HHstar(U+V) \leq \HHstar(U)\HHstar(V)
\end{equation}
and
\begin{equation}
\label{eq: Schmidt's inequality L_xi}
L_\xi(U \cap V) L_\xi(U + V) \leq L_\xi(U)L_\xi(V).
\end{equation}
\end{Prop}

\begin{proof}
In view of \eqref{eq: D_xi(V) = L_xi(V^perp)}, we only need to prove that \eqref{eq: Schmidt's inequality L_xi} holds for any pair $(U,V)$ as in the statement of the proposition (for then, it suffices to apply \eqref{eq: Schmidt's inequality L_xi} to the pair $(U^\perp,V^\perp)$). We follow Schmidt's arguments \cite[Ch.\,I, Lem.\,8A]{Schmidt}. For any pure products $\bX, \bY, \bZZ \in \sbw{\,}{\bR^{m+1}}$, we have
\begin{equation}
\label{eq proof: prop Schmidt inegalite pour HHstar}
\norm{\bX} \norm{\bX\wedge \bY \wedge \bZZ} \leq \norm{\bX\wedge \bY} \norm{\bX\wedge \bZZ}.
\end{equation}
Let $U, V$ be subspaces of $\bR^{m+1}$ defined over $\bQ$. If $U=\{0\}$ or $V = \{0\}$, then \eqref{eq: Schmidt's inequality L_xi} is trivial, so we may assume that $U$ and $V$ have dimension $\geq 1$. Let $\bx_1,\dots,\bx_r$ be a $\bZ$-basis of $U\cap V \cap \bZ^{m+1}$, which we complete to a $\bZ$-basis $\bx_1,\dots,\bx_r,\by_1,\dots,\by_s$ of $U\cap\bZ^{m+1}$ (\resp $\bx_1,\dots,\bx_r,\bz_1,\dots,\bz_t$ of $V\cap\bZ^{m+1}$). Set
\begin{equation*}
\bX = \Xi_m\wedge \bx_1 \wedge \dots \bx_r, \quad \bY = \by_1\wedge \dots \wedge \by_s \AND \bZZ = \bz_1 \wedge \cdots \wedge\bz_t.
\end{equation*}
We get \eqref{eq: Schmidt's inequality L_xi} by applying \eqref{eq proof: prop Schmidt inegalite pour HHstar} to the above pure products.
\end{proof}

We identify $\bR[X]_{\leq m}$ to $\bR^{m+1}$ and $\bR^{m+1}$ to the space of $(m+1)\times 1$ column matrices with real coefficients via the isomorphism
\begin{equation}
\label{eq: identification poly avec points de R^m}
\sum_{k=0}^{m} a_k X^k \longmapsto (a_0,\dots,a_{m}) \AND
(a_0,\dots,a_{m})\mapsto
\begin{pmatrix}
a_0 \\
\vdots \\
a_{m}
\end{pmatrix}.
\end{equation}
Then, for any $P\in\bR[X]_{\leq m}$, we have $P(\xi) = \psc{\bz}{\Xi_m}$, where $\bz\in\bR^{m+1}$ corresponds to $P$. In particular, if $P\in\bZ[X]_{\leq m}$ is primitive, then \eqref{eq: HHstar(bx)} may be rewritten as
\begin{equation}
\label{eq: HHstar(P)}
\HHstar\big(\Vect[\bR]{P}\big) = |P(\xi)|.
\end{equation}
We will repeatedly use the following ``twisted'' dual version of \cite[Lem.\,2.1]{poels2022simultaneous}.

\begin{Lem}
\label{lem: estimation des pseudo resultant tordus}
There is a positive constant $c$, which only depends on $n$ and $\xi$, with the following property. For any linearly independent polynomials $P_1,\dots,P_k\in\bZ[X]_{\leq m}$ (with $k\geq 1$), we have
\begin{equation}
\label{eq: estimation L_xi(V) avec T}
\HHstar\big(\Vect[\bR]{P_1,\dots,P_k} \big) \leq c \sum_{i=1}^{k} \frac{|P_i(\xi)|}{\normH{P_i}} \prod_{j=1}^{k} \normH{P_j}.
\end{equation}
\end{Lem}
Note that for any $P\in\bZ[X]_{\leq m}$, Eq.~\eqref{eq: HHstar(P)} implies that $\HHstar(\Vect[\bR]{P}) \leq |P(\xi)|$.

\begin{proof}
Let $Q_1,\dots,Q_k$ be a $\bZ$-basis of $V\cap\bZ[X]_{\leq m}$, where $V=\Vect[\bR]{P_1,\dots,P_k}$. There exists a non-zero $\alpha\in\bZ$ such that
\[
P_1\wedge \cdots \wedge P_k = \alpha Q_1\wedge \cdots \wedge Q_k,
\]
and so
\[
\HHstar(V) = \norm{ \Xi_m\,\lrcorner \,(Q_1\wedge \cdots \wedge Q_k)} \leq \norm{ \Xi_m\,\lrcorner \,(P_1\wedge \cdots \wedge P_k)}.
\]
On the other hand, by \eqref{eq: contraction formule explicite} combined with Hadamard's inequality, we obtain
\begin{align*}
\norm{\Xi_m\,\lrcorner \,(P_1\wedge \cdots \wedge P_k)} & = \bnorm{\sum_{i=1}^{k} (-1)^{k-i} P_i(\xi) \cdot P_1\wedge \cdots \wedge \widehat{P_i} \wedge \cdots \wedge P_k } \\
& \ll \sum_{i=1}^{k} |P_i(\xi)| \normH{P_1}\cdots \widehat{\normH{P_i}} \cdots \normH{P_k}
\end{align*}
(recall that the norm $\normH{\cdot}$ is defined in Section~\ref{section: notation}).
\end{proof}

\section{Subfamilies of polynomials: dimension and height}
\label{section : indpce lineaire Q_i}

Let $d,n,\xi$ and $\homega$ be as in Section~\ref{section: sequence des Q_i}. In particular we have
\[
2\leq d < 1+\frac{n}{2},
\]
and we suppose that \eqref{eq: notation homega} holds, namely
\[
\homega_n(\xi) > \homega > 2n -d.
\]
Let us fix a sequence of minimal polynomials $(P_i)_{i\geq 0}$ associated to $n$ and $\xi$ as in Section~\ref{section: minimal pol}. We denote by $(Q_i)_{i\geq 0}$ the sequence of irreducible factors given by Proposition~\ref{prop: existence Qi}. In particular, for each $i\geq 0$ we have
\begin{equation}
\label{eq: petitesse de Q_i(xi)}
|Q_i(\xi)| \leq \normH{Q_i}^{-\homega},
\end{equation}
as well as
\begin{equation}
\label{eq: taille de H(Q_i+1)}
\normH{Q_{i+1}}^\tau \leq \normH{Q_{i}}, \quad \textrm{where } \tau = \frac{\homega\big(\homega -n-d+3\big)}{\omega_n(\xi)\big(\omega_n(\xi)-n-d+3\big)} \in [0,1).
\end{equation}
Assuming that $d$ is not too large, we will prove in the next section that $\omega_n(\xi) < \infty$, and thus $\tau > 0$.
Here, we investigate the following question: can we find ``large'' subfamilies of $(Q_i)_{i\geq 0}$ which are linearly independent, and whose elements have ``comparable'' height? More precisely, given two indices $k< i$, can we find an exponent $\theta_j\in(0,1)$ which depends only on $d,n$ and the dimension $j+1$ of the subspace $\Vect[\bR]{Q_k,Q_{k+1},\dots,Q_i}$ (and not on the indices $i$ and $k$), such that $\normH{Q_i}^{\theta_j} \ll \normH{Q_k}$? For $i=k+1$, we already have \eqref{eq: taille de H(Q_i+1)}. With this goal in mind, let us introduce some notation.

\begin{Def}
\label{Def m_n}
Let $m_n = m_n(\xi) \in [2,n+1]$ be the integer
\[
m_n := \lim_{i\to\infty} \dim(\Vect[\bR]{Q_i,Q_{i+1},\dots}).
\]
\end{Def}

\begin{remark*}
Note that we might have $m_n < n+1$, since, unlike for simultaneous approximation (see \cite[Eq.\,(5.3)]{poels2022simultaneous}), it is possible that the sequence $(P_i)_{i\geq j}$ is contained in a proper subspace of $\bR[X]_{\leq n}$, see \eg \cite{moshchevitin2007best}. However, we will show later that under the hypothesis $d \asymp n^{1/3}$, we have $m_n\gg n^{1/3}$. The next definition is somewhat dual to \cite[Def.\,5.2]{poels2022simultaneous}. However, note that in \cite[Def.\,5.2]{poels2022simultaneous}, the sets $A_j[i]$ are constructed from the points $\bx_i,\bx_{i+1},\dots$ coming \emph{after} the good approximation $\bx_i$, whereas in the present setting we need to consider the points $Q_{i}, Q_{i-1},\dots$ coming \emph{before} $Q_i$. It does not seem to work well the other way round.
\end{remark*}

\begin{Def}
\label{Def: j_0,j_1,sigma_j(i) et Aj[i]}
Let $j_1 > j_0\geq 0$ be such that
\[
\dim \Vect[\bR]{Q_{j_0},Q_{j_0+1},\dots, Q_{j_1}} = \dim \Vect[\bR]{Q_{j_0},Q_{j_0+1},\dots} = m_n.
\]
For each $i\geq j_1$ and $j=0,\dots, m_n -2$, we define
\[
\sigma_j(i) = k, \quad A_j[i] = \{Q_k, Q_{k+1},\dots, Q_i \} \AND Y_j(i) = \normH{Q_{k-1}},
\]
where $k\in\{j_0+1,\dots,i\}$ is the smallest index such that $\dim \Vect[\bR]{Q_k,\dots, Q_i} = j+1$.
\end{Def}

Proposition~\ref{Cor: V_2n-k = tout l'espace} implies that
\begin{equation}
\label{eq: V_2n-j(A_j[i]) = tout l'espace}
V_{2n-j}\big(A_j[i]\big) = \bR[X]_{\leq 2n-j} \qquad (j=1,\dots,m_n-2).
\end{equation}

\begin{Def}
\label{Def: tau_j}
Let $\tau\in(0,1)$. We associate to $\tau$ a sequence $(\tau_j)_{0\leq j\leq n/2}$ by setting $\tau_0 = \tau$, and for $j=1,\dots,\lfloor n/2 \rfloor$
\begin{equation*}
\tau_{j} = \alpha_j\Bigl(\tau_{j-1} - \frac{2j-1}{2n-d} \Bigr), \quad \textrm{where } \alpha_j = \frac{(2n-d)\tau^2}{(n-2j)\tau+n-j+1}.
\end{equation*}
\end{Def}

The first main result of this section is the following.

\begin{Prop}
\label{Prop: prop resumee avec tau_j}
Let $\tau\in(0,1)$ and let $(\tau_j)_{0\leq j \leq n/2}$ be as in Definition~\ref{Def: tau_j}. Suppose that
\begin{equation}
\label{eq: def exposant tau}
\normH{Q_{i+1}}^\tau \leq \normH{Q_i} \qquad \textrm{for each sufficiently large $i$}.
\end{equation}
Then for each large enough $i$, we also have
\begin{equation}
\label{eq Prop: minoration de Y_j}
\normH{Q_i}^{\tau_j} \ll Y_j(i) \qquad \textrm{for } j=0,\dots, \min\big\{\lfloor n/2\rfloor,m_n-2\big\},
\end{equation}
with implicit constants which do not depend on $i$ and $j$.
\end{Prop}

\begin{remark*}
We will use the exponent $\tau$ given by \eqref{eq: taille de H(Q_i+1)}. We will prove that under suitable conditions, the exponent of best approximation $\omega_n(\xi)$ is not ``too large''. This will ensure that $\tau$ is ``close'' to~$1$. This issue, which is one of the delicate parts of this paper, will be dealt with in Section~\ref{section: estimation de omega_n}.
\end{remark*}

In order to get \eqref{eq Prop: minoration de Y_j}, we will try to adopt a strategy similar to the one of \hbox{\cite[\S 5]{poels2022simultaneous}} in the setting of simultaneous approximation to the successive powers of $\xi$. New difficulties arise however, for example we need to work with $\HHstar$ instead of the standard height of subspaces (see Section~\ref{section: hauteurs tordues}). Schmidt's inequality \eqref{eq: Schmidt's inequality HHstar} will play a key-role in our proofs. We use the notation of Definition~\ref{Def: fonction g_A(N)} for the sets $\cB_k(\cA)$ and the subspaces $V_k(\cA) \subset \bR[X]_{\leq k}$.

\begin{proof}
Without loss of generality, we may suppose that the index $j_0$ is large enough so that \eqref{eq: def exposant tau} holds for each $i\geq j_0-1$. Fix $i\geq j_1$, and for simplicity write $m=m_n$ and $Y_k:= Y_k(i)$ for $k=0,\dots,m-2$.

We prove \eqref{eq Prop: minoration de Y_j} by induction on $j$. If $j=0$, we have $Y_0 = \normH{Q_{i-1}}$ since $\sigma_0(i) = i$. By~\eqref{eq: def exposant tau} applied with $i-1$ instead of $i$, we get $\normH{Q_i}^{\tau_0} \leq Y_0$. Now, let $j\in\{1,\dots,m-2\}$ with $j\leq n/2$ such that \eqref{eq Prop: minoration de Y_j} holds for $j-1$. If $\tau_{j}\leq 0$, then \eqref{eq Prop: minoration de Y_j} holds trivially for $j$. We assume that $\tau_{j} > 0$. Consequently, we also have $\tau_{j-1} > 0$. Write $P:= Q_{\sigma_j(i)}$ and $Q:= Q_{\sigma_j(i)+1}$. By \eqref{eq: def exposant tau}, we have
\begin{equation}
\label{eq proof: contrôle H(P) et H(Q) par Y_j}
\normH{Q}^{\tau^2} \leq \normH{P}^{\tau} \leq Y_j.
\end{equation}
Since $P$ and $Q$ are coprime, Lemma~\ref{lem : dim V_n+j(P,Q)} implies that $\dim V_{2n-j}(P,Q)\geq 2(n-j+1)$. Therefore, there exists a family of $2n-3j+1$ linearly independent polynomials
\[
\cU_j:=\{U_0,\dots,U_{2n-3j}\} \subset \cB_{2n-j}(P,Q)
\]
such that $\Vect[\bR]{A_j[i]} \cap \Vect[\bR]{\cU_j} = \{0\}$. Note that since $j\leq n/2$, we may choose $\cU_j$ such that it contains at least $n-2j$ polynomials whose height is equal to $\normH{P}$. The remaining $n-j+1$ ones have height~$\leq \normH{Q}$. By \eqref{eq: V_2n-j(A_j[i]) = tout l'espace}, we have $V_{2n-j}\big(A_j[i]\big) = \bR[X]_{\leq 2n-j}$. Therefore, there exists
\[
\cV_j:=\{V_1,\dots,V_{j-1}\} \subset \cB_{2n-j}(A_j[i]) = \cB_{2n-j}(Q_{\sigma_j(i)},\dots, Q_i)
\]
(with the convention $\cV_j = \emptyset$ if $j=1$) such that we have the direct sum
\[
\Vect[\bR]{A_j[i]} \oplus \Vect[\bR]{\cU_j} \oplus \Vect[\bR]{\cV_j} = \bR[X]_{\leq 2n-j}.
\]
All the polynomials of $\cV_j$ have height at most $\normH{Q_i}$. Let $k\in\{\sigma_j(i),\dots,i\}$ which maximizes $|Q_k(\xi)|/\normH{Q_k}$ and define
\[
A:= \Vect[\bR]{A_j[i]} \AND B:= \Vect[\bR]{\cU_j \cup \cV_j \cup \{Q_k\}},
\]
so that $A + B = \bR[X]_{\leq 2n-j}$ and $A\cap B = \Vect[\bR]{Q_k}$. We will now make a crucial use of the function $\HHstar$ introduced in Definition~\ref{Def: def HHstar} (here, the ambient space is $\bR[X]_{\leq 2n-j}$, identified to $\bR^{2n-j+1}$ via \eqref{eq: identification poly avec points de R^m}). Recall that
\[
\HHstar(A+B) = \HHstar\big(\bR[X]_{\leq 2n-j}\big) = \norm{(1,\xi,\dots,\xi^{2n-j})} \asymp 1,
\]
and that according to \eqref{eq: HHstar(P)} the primitive polynomial $Q_k$ satisfies
\[
\HHstar(A\cap B) = \HHstar\big(\Vect[\bR]{Q_k}\big) = |Q_k(\xi)|.
\]
Schmidt's inequality \eqref{eq: Schmidt's inequality HHstar} applied with the subspaces $A$ and $B$ yields
\begin{equation}
\label{eq proof: Schmidt notre construction}
|Q_{k}(\xi)| \asymp \HHstar(A+B)\HHstar(A\cap B) \leq \HHstar(A) \HHstar(B),
\end{equation}
the implicit constants depending only on $n$ and $\xi$ (and not on the indices $i,j$). It~\hbox{remains} to estimate $\HHstar(A)$ and $\HHstar(B)$. The subspace $B\subset\bR[X]_{\leq 2n-j}$ is generated by the $2n-2j+1$ linearly independent polynomials $\cV = \cU_j \cup \cV_j \cup \{Q_k\}$. Moreover (see the remarks after the constructions of $\cU_j$ and $\cV_j$), we have
\[
\prod_{R\in\cV}\normH{R} \leq \normH{P}^{n-2j}\normH{Q}^{n-j+1}\normH{Q_i}^{j-1}\normH{Q_k}.
\]
By choice of $k$, for each $R\in\cV$ we also have $|R(\xi)|/\normH{R}\ll |Q_k(\xi)|/\normH{Q_k}$, and Lemma~\ref{lem: estimation des pseudo resultant tordus} combined with the above yields the upper bound
\begin{equation*}
\HHstar(B) \ll |Q_k(\xi)| \normH{P}^{n-2j}\normH{Q}^{n-j+1}\normH{Q_i}^{j-1}.
\end{equation*}
The space $A = \Vect[\bR]{A_j[i]}\subset\bR[X]_{\leq 2n-j}$ is spanned by a set $\cU$ of $j+1$ linearly polynomials that may be chosen among $Q_{\sigma_{j-1}(i)-1}$,...,$Q_{i-1}$, $Q_i$. For each $R\in\cU$, we have $\normH{R}\leq \normH{Q_i}$ and $|R(\xi)| \leq \normH{R}^{-\homega} \leq Y_{j-1}^{-\homega}$. Combined with Lemma~\ref{lem: estimation des pseudo resultant tordus}, we obtain
\begin{equation*}
\HHstar(A) \ll \sum_{R\in\cU}|R(\xi)|\prod_{\substack{S\in\cU \\ S\neq R}}\normH{S} \ll Y_{j-1}^{-\homega}\normH{Q_{i}}^j.
\end{equation*}
Then, combining the above upper bounds for $\HHstar(B)$ and $\HHstar(A)$ with \eqref{eq proof: Schmidt notre construction} and \eqref{eq proof: contrôle H(P) et H(Q) par Y_j}, we get
\begin{equation*}
Y_{j-1}^{\homega} \ll \normH{P}^{n-2j}\normH{Q}^{n-j+1}\normH{Q_i}^{2j-1} \ll Y_j^{(n-2j)/\tau + (n-j+1)/\tau^2}\normH{Q_i}^{2j-1},
\end{equation*}
where the implicit constants depend on $n$ and $\xi$ only. Using the induction hypothesis, we also have $\normH{Q_i}^{\homega\,\tau_{j-1}} \ll Y_{j-1}^{\homega}$, hence
\begin{equation*}
\normH{Q_i}^{\homega\,\tau_{j-1}-2j+1} \ll Y_j^{(n-2j)/\tau + (n-j+1)/\tau^2} = Y_j^{(2n-d)/\alpha_j}.
\end{equation*}
Rising each term to the power $\alpha_j/(2n-d)$ and using $\homega > 2n-d$, we easily deduce \eqref{eq Prop: minoration de Y_j} for $j$. This concludes our induction step.
\end{proof}

\begin{Rem}
We could get a slightly greater exponent $\tau_j$ in the above proposition by using a more precise estimate for $\HHstar(A)$. However, this improvement would at best lead to a larger constant $a$ in Theorem~\ref{Thm : main} ; the term $n^{1/3}$ would remain the same, whereas we are expecting $n^{1/2}$. We preferred to keep the arguments simple.
\end{Rem}

The following result is inspired by Laurent's approach in \cite[Lem.\,5]{laurent2003simultaneous}.

\begin{Prop}
\label{Prop: majoration Y_m-2}
Let the hypotheses be as in Proposition~\ref{Prop: prop resumee avec tau_j} and write $m=m_n$. For any $\lambda < \lambda_n(\xi)$, there are infinitely many indices $i$ such that
\begin{equation*}
\label{eq: majoration theta_m-2 eq 0}
Y_{m-2}(i) < \normH{Q_i}^{1/(\homega\lambda\tau)}.
\end{equation*}
In particular, there are infinitely many indices $i$ such that
\begin{equation}
\label{eq: majoration theta_m-2}
Y_{m-2}(i) \leq \normH{Q_i}^\mu,\quad \textrm{where } \mu := \frac{n}{(2n-d)\tau}.
\end{equation}
\end{Prop}

\begin{proof}
By definition of $m$, the subspace
\begin{equation}
\label{eq proof: def V}
V = \Vect[\bR]{Q_{\sigma_{m-2}(i)-1},Q_{\sigma_{m-2}(i)},\dots,Q_{i}}
\end{equation}
of $\bR[X]_{\leq n}$ is independent of $i$ for $i\geq j_1$, where $j_1$ comes from Definition~\ref{Def: j_0,j_1,sigma_j(i) et Aj[i]}. It has dimension $m$ since $\dim A_{m-2}[i] = m-1$ and $Q_{\sigma_{m-2}(i)-1}\notin A_{m-2}[i]$. Fix two positive real numbers $\alpha,\lambda$ with $\lambda < \alpha < \lambda_n(\xi)$, and suppose by contradiction that there exists an index $i_0\geq j_1$ such that for each $i \geq i_0$
\begin{equation}
\label{eq proof: def theta}
Y_{m-2}(i) \geq \normH{Q_i}^\theta, \quad \textrm{where } \theta=\frac{1}{\homega\lambda\tau}.
\end{equation}
By hypothesis, we can also assume that $\normH{Q_{i+1}}^\tau \leq \normH{Q_i}$ for each $i\geq i_0$. Identifying $\bR[X]_{\leq n}$ with $\bR^{n+1}$ via the isomorphism \eqref{eq: identification poly avec points de R^m}, we claim that the point $\Xi=(1,\xi,\xi^2,\dots,\xi^n)$ is orthogonal to $V$, with respect to the standard scalar product $\psc{\cdot}{\cdot}$ of $\bR^{n+1}$.

By definition of $\lambda_n(\xi)$, there exist infinitely many non-zero $\by=(y_0,\dots,y_n)\in\bZ^{n+1}$ satisfying
\begin{equation*}
L(\by) =\max_{1\le k \le n} |y_0\xi^k - y_k| \leq Y^{-\alpha},\quad \textrm{where } Y = \normH{\by} = \max_{1\leq k\leq n} |y_k|.
\end{equation*}
Let $(\by_i)_{i\geq 0}$ be an unbounded sequence of such points ordered by increasing norm. This sequence converges projectively to $\Xi = (1,\xi,\xi^2,\dots,\xi^n)$. Without loss of generality, we may assume that $\normH{\by_0}^{\alpha} > 2(n+1)\normH{Q_{i_0}}$. Fix an index $j$ arbitrarily large. For simplicity, set $\by:= \by_j$ and $Y = \normH{\by_j}$. There exists an index $i\geq i_0$ such that
\begin{equation}
\label{eq proof: sur l'exposant mu eq1}
\normH{Q_i} < \frac{Y^{\alpha}}{2(n+1)} \leq \normH{Q_{i+1}} \leq \normH{Q_i}^{1/\tau}.
\end{equation}
Note that $i$ tends to infinity with $j$. Let $k\in\{\sigma_{m-2}(i)-1,\dots,i\}$. The polynomial $Q:=Q_k$ is identified with an integer point $\bz\in\bZ^{n+1}$ such that $Q(\xi)= \psc{\bz}{\Xi}$. Since $\psc{\bz}{\by} = \psc{\bz}{\by -y_0\Xi} + y_0\psc{\bz}{\Xi}$, we get
\begin{equation*}
|\psc{\bz}{\by}| \leq (n+1)\normH{Q}L(\by) + Y |Q(\xi)|
\end{equation*}
(cf. \cite[Lem.\,5]{laurent2003simultaneous}). Our hypothesis \eqref{eq proof: def theta} yields
\begin{equation*}
\normH{Q_i}^\theta \leq Y_{m-2}(i) \leq \normH{Q} \leq \normH{Q_i}.
\end{equation*}
Using \eqref{eq proof: sur l'exposant mu eq1} together with $L(\by)\leq Y^{-\alpha}$, we get
\begin{equation*}
(n+1)\normH{Q}L(\by) <\frac{1}{2}.
\end{equation*}
Moreover, \eqref{eq proof: sur l'exposant mu eq1} also yields $ Y^{1/\homega} \ll \normH{Q_i}^{1/(\homega\alpha\tau)} = \normH{Q_i}^{\theta \lambda / \alpha}$, where the implicit constant only depends on $n$. Since $\lambda < \alpha$, we may choose $j$ so large that $(2Y)^{1/\homega} < \normH{Q_i}^\theta$. Combining this with the estimate $|Q(\xi)|\leq \normH{Q}^{-\homega}$ from \eqref{eq: petitesse de Q_i(xi)}, we also get
\[
Y|Q(\xi)|\leq Y\normH{Q}^{-\homega} \leq Y\normH{Q_i}^{-\theta\,\homega} < \frac{1}{2}.
\]
We conclude that the integer $|\psc{\bz}{\by}|$ is (strictly) less that $1$. It is thus equal to $0$, and so $\by$ and $\bz$ are orthogonal. By letting $k$ vary, this implies that $\by = \by_j$ is orthogonal to the subspace $V$. Since this is true for all sufficiently large $j$, it follows that the (projective) limit $\Xi$ is also orthogonal to $V$. This proves our claim and provides the required contradiction since no $Q_i$ vanishes at the transcendental number $\xi$. Thus, \eqref{eq proof: def theta} does not hold for arbitrarily large indices $i$. Estimate \eqref{eq: majoration theta_m-2} follows by noticing that $\lambda_n(\xi)\geq 1/n$ by Dirichlet's theorem, and recalling that $\homega > 2n-d$. We may therefore choose $\lambda < \lambda_n(\xi)$ so that $\lambda\homega > (2n-d)/n$.
\end{proof}

\begin{Cor}
\label{Cor: tau_m_2 < mu}
Under the same hypotheses, suppose moreover that $m=m_n$ satisfies $m-2\leq n/2$, and let $(\tau_j)_{0\leq j \leq n/2}$ be as in Definition~\ref{Def: tau_j}. Then, we have
\[
\tau_{m-2} \leq \mu = \frac{n}{(2n-d)\tau}.
\]
\end{Cor}

\begin{proof}
By Propositions~\ref{Prop: prop resumee avec tau_j} and~\ref{Prop: majoration Y_m-2} there are infinitely many indices $i$ for which $\normH{Q_i}^{\tau_{m-2}} \ll Y_{m-2}(i) \leq \normH{Q_i}^\mu$. Since $\normH{Q_i}$ tends to infinity with $i$, we deduce that $\tau_{m-2}\leq \mu$.
\end{proof}

\section{Upper bound on the exponent of best approximation}
\label{section: estimation de omega_n}

This section is devoted to the proof of the following upper bound for $\omega_n(\xi)$.

\begin{Prop}
\label{Prop: estimation omega_n(xi)}
Suppose that $\homega_n(\xi) > 2n-d$, with an integer $d\in\bN$ satisfying $\displaystyle 2 \leq d\leq \sqrt[3]{n/4}$. Then, we have the upper bound
\begin{equation*}
\omega_n(\xi)\leq 2n + P(n,d), \quad\textrm{where } P(n,d) = \frac{n(4d^2-d-5)+8d^2-2d-15}{2n-8d^2+2d+15}.
\end{equation*}
If moreover we have $d\leq \Big\lceil\sqrt[3]{n/16} \Big\rceil$ and $n > 16$, then
\begin{equation*}
\omega_n(\xi)\leq 2n + 2d^2.
\end{equation*}
\end{Prop}

Let $d,n,\xi$ and $\homega$ be as in Sections~\ref{section: sequence des Q_i} and~\ref{section : indpce lineaire Q_i}. We suppose thus that $2\leq d < 1+n/2$ and that \eqref{eq: notation homega} holds, namely
\[
\homega_n(\xi) > \homega > 2n -d.
\]
Fix a sequence of minimal polynomials $(P_i)_{i\geq 0}$ associated to $n$ and $\xi$ as in Section~\ref{section: minimal pol}. We denote by $(Q_i)_{i\geq 0}$ the sequence of irreducible factors given by Proposition~\ref{prop: existence Qi}. Unless otherwise stated, all the constants implicit in the symbols $\ll$, $\gg, \asymp$ and $\GrO(\cdot)$ depend only on $n$, $d$, $\xi$ and $\homega$.

According to Proposition~\ref{prop: existence Qi}, we have $\omega_n(\xi) = \limsup_{i\to\infty}\omega(Q_i)$. By \eqref{eq: premiere estimation Pj vs Qi}, we also have
\begin{equation}
\label{eq: rappel theta_i}
\normH{P_j} \leq \normH{Q_i}^{1+\theta_i} \quad \textrm{with } \theta_i = \frac{\omega(Q_i)-2n+d}{n-2d+3},
\end{equation}
for each $i\geq 0$ and each $j$ such that $Q_i$ divides $P_j$. Proposition~\ref{Prop: estimation omega_n(xi)} implies that if $d^3$ is small compared to $n$, then $\theta_i = \GrO(d^2/n)$ is small, and $Q_i$ has ``almost'' the same norm as $P_j$.

In order to bound from above $\omega_n(\xi)$, it suffices to do so for $\omega(Q_i)$. We could try to use \eqref{eq: rappel theta_i}, which implies that any minimal polynomial of height greater than $\normH{Q_i}^{1+\theta_i}$ is not divisible by $Q_i$. They are thus coprime and we may consider their (non-zero) resultant. However we cannot conclude, as $\theta_i$ is too large. To resolve this problem, we~need several lemmas. We first start by a few simple observations. A quick computation yields
\begin{equation}
\label{equation fonctionnelle theta eq 1}
(1+\theta_i)(2n-d) = \omega(Q_i) + (n+d-3)\theta_i.
\end{equation}
More generally, for each $\eta\geq 0$, we have
\begin{equation}
\label{equation fonctionnelle theta eq 2}
\big(1+\theta_i(1-\eta) \big)(2n-d) = \omega(Q_i) + \big( n+d-3 - \eta(2n-d) \big)\theta_i.
\end{equation}
Under the condition $\eta < (n+d-3)/(2n-d)$, which holds as soon as $\eta < 1/2$, this implies that for each $i\geq 0$, we have
\begin{equation}
\label{equation: minoration de Q_i(xi) via theta_i}
|Q_i(\xi)| = \normH{Q_i}^{-\omega(Q_i)} > \normH{Q_i}^{-\big(1+\theta_i(1-\eta)\big)(2n-d)}.
\end{equation}

\begin{Lem}
\label{Lem: P = QR grand bonne approx}
Let $i\geq 0$ and $\eta\in[0,1/2)$, and suppose that $R\in\bZ[X]_{\leq d-2}$ is a non-zero polynomial such that $P := Q_iR$ has degree at most $n$ and satisfies
\begin{equation}
\label{eq : systeme dio pour P = QR}
\normH{P} \leq H:= \normH{Q_i}^{1+\theta_i(1-\eta)} \AND |P(\xi)| \leq H^{-2n+d}.
\end{equation}
Define
\begin{equation*}
\eta' = \frac{(2n-d)\eta}{n+d-3} \AND \eta'' = \frac{(2n-2d+3)\eta+d-3}{n+d-3}.
\end{equation*}
Then, we have the following properties.
\begin{enumerate}
\item \label{item: item 1 Lem P = QR sol }
The polynomial $R$ is non-constant. We have $d\geq 3$ and
\begin{equation}
\label{eq Prop: P = QR grand bonne approx: R(xi) revisite}
\normH{R}^{-(n+d-3)} \ll |R(\xi)| \leq \normH{Q}^{-(n+d-3)(1-\eta')\theta_i}.
\end{equation}
\item \label{item: item 2 Lem P = QR sol }
There exist a non-constant irreducible polynomial $A\in\bZ[X]_{\leq n}$ and an integer $e \in[1,d-2]$ such that $A^e$ divides $R$,
\begin{equation}
\label{eq Prop: P = QR grand bonne approx: H(A^e) revisite}
\normH{A^e} \gg \normH{Q}^{\theta(1-\eta'')} \AND \normH{A^e}^{-(n+d-3)} \ll |A^e(\xi)|.
\end{equation}
\item Let $A$ and $e$ be as in \eqref{item: item 2 Lem P = QR sol }. If $S\in\bZ[X]_{\leq d-2}$ is a non-zero polynomial such that $A$ and $S$ are coprime and $\normH{S}\leq \normH{A^e}$, then
\begin{equation}
\label{eq Prop: P = QR grand bonne approx: H(S) revisite}
|S(\xi)| \gg \normH{A^e}^{-(2d-5)}.
\end{equation}
\end{enumerate}
\end{Lem}

\begin{proof}
Fix $i\geq 0$. For simplicity, write $Q:= Q_i$ and $\theta = \theta_i$. By Gelfond's lemma, we have
\begin{equation*}
\normH{Q}\normH{R}\asymp \normH{QR} = \normH{P} \leq \normH{Q}^{1+\theta(1-\eta)},
\end{equation*}
so that
\begin{equation}
\label{eq proof: H(R) << fonction de H(Q)}
\normH{R}\ll \normH{Q}^{\theta(1-\eta)}.
\end{equation}
The first inequality of \eqref{eq Prop: P = QR grand bonne approx: R(xi) revisite} and the second of \eqref{eq Prop: P = QR grand bonne approx: H(A^e) revisite} are consequences of Lemma~\ref{lemme: omega(R) petit degre} (using $\deg (A^e) \leq \deg(R) \leq d-2$). Using \eqref{equation fonctionnelle theta eq 2} together with \eqref{eq : systeme dio pour P = QR} and $\normH{Q}^{-\omega(Q)} = |Q(\xi)|$, we find
\begin{equation*}
|Q(\xi)R(\xi)| = |P(\xi)| \leq \normH{Q}^{-\big(1+\theta(1-\eta)\big)(2n-d)} = |Q(\xi)|\normH{Q}^{- \big( n+d-3 - \eta(2n-d) \big)\theta}.
\end{equation*}
Simplifying by $|Q(\xi)|$ yields the second inequality of \eqref{eq Prop: P = QR grand bonne approx: R(xi) revisite}. In particular we have $|R(\xi)| < 1$ since $\normH{Q} > 1$ (and $\theta > 0$ as well as $\eta'\leq 2\eta < 1$). Consequently $R\in\bZ[X]_{\leq d-2}$ cannot be constant, and thus $d\geq 3$.

Without loss of generality, we may suppose that $P$ (and thus $R$) is primitive. Let us consider the factorization of $R$ over $\bZ$. There exist an integer $k\geq 1$, irreducible (non-constant) pairwise distinct polynomials $A_1,\dots,A_k\in\bZ[X]$ and positive integers $\alpha_1,\dots,\alpha_k$ such that
\[
R = \prod_{j=1}^{k} A_j^{\alpha_j} = \prod_{j=1}^{k} B_j \qquad \textrm{with $B_j:= A_j^{\alpha_j}$ for each $j=1,\dots,k$}.
\]
According to Lemma~\ref{lem: lem general 2}, there exists $j\in\{1,\dots,k\}$ such that $B = B_j$ satisfies
\begin{equation*}
|B(\xi)| \ll \normH{R}^{d-3}|R(\xi)|.
\end{equation*}
Set $A=A_j$ and $e=\alpha_j$. We use \eqref{eq proof: H(R) << fonction de H(Q)} to bound $\normH{R}$ from above, and the second inequality of~\eqref{eq Prop: P = QR grand bonne approx: R(xi) revisite} to bound $|R(\xi)|$ from above. Then, Lemma~\ref{lemme: omega(R) petit degre} applied to the polynomial $B\in\bZ[X]_{\leq d-2}$ together with the above yields
\begin{equation*}
\normH{B}^{-(n+d-3)} \ll |B(\xi)| \ll \normH{R}^{d-3}|R(\xi)| \ll \normH{Q}^{(d-3)(1-\eta)\theta-(n+d-3)(1-\eta')\theta},
\end{equation*}
Since by definition of $\eta'$ and $\eta''$ we have
\begin{equation*}
1-\eta'-\frac{d-3}{n+d-3}(1-\eta) = 1-\eta'',
\end{equation*}
we deduce that
\begin{equation}
\label{eq proof: B(xi) << fonction de H(Q)}
\normH{B}^{-(n+d-3)} \ll |B(\xi)| \ll \normH{Q}^{-(n+d-3)\theta(1-\eta'')}.
\end{equation}
and \eqref{eq Prop: P = QR grand bonne approx: H(A^e) revisite} follows easily upon recalling that $A^e = B$. Now, suppose that $S\in\bZ[X]_{\leq d-2}$ is a non-zero polynomial coprime to $A$ with $\normH{S}\leq \normH{B}$. If $S$ is constant, then \eqref{eq Prop: P = QR grand bonne approx: H(S) revisite} is trivial. We may therefore assume that $S$ has degree at least $1$. Then, the estimate of Lemma~\ref{lem: estimation classique du resultant} yields
\begin{equation}
\begin{aligned}
1 \leq |\Res(B,S)| & \ll \normH{B}^{d-3}\normH{S}^{d-2}|B(\xi)| + \normH{B}^{d-2}\normH{S}^{d-3}|S(\xi)| \\
& \ll \normH{B}^{2d-5}\big(|B(\xi)|+|S(\xi)| \big) \label{eq proof: estimation resultant}
\end{aligned}
\end{equation}
(where the implicit constants depend on $\xi$, $n$ and $c$). As $B$ divides $R$, we have $\normH{B}\ll \normH{R}$. Together with \eqref{eq proof: H(R) << fonction de H(Q)}, this gives $\normH{B} \ll \normH{Q}^{\theta(1-\eta)}$. Combining the above with \eqref{eq proof: B(xi) << fonction de H(Q)}, we obtain
\begin{equation*}
\normH{B}^{2d-5}|B(\xi)| \ll \normH{Q}^{(2d-5)\theta(1-\eta)-(n+d-3)\theta(1-\eta'')}.
\end{equation*}
On the other hand, using $\eta\leq 1/2$ we get
\begin{equation*}
(2d-5)(1-\eta)-(n+d-3)(1-\eta'')= (2n-4d+8)\eta - (n-2d+5) \leq -1.
\end{equation*}
Since for each large enough $i$, the number $\theta = \theta_i$ is bounded from below by
\[
\rho = \frac{\homega-2n+d}{n-2d+3} > 0,
\]
it follows that $ \normH{B}^{2d-5}|B(\xi)| \ll \normH{Q}^{-\rho}$ tends to $0$ as $i$ tends to infinity. Consequently, \eqref{eq proof: estimation resultant} becomes
\begin{equation*}
1 \ll \normH{B}^{2d-5}|S(\xi)|,
\end{equation*}
hence \eqref{eq Prop: P = QR grand bonne approx: H(S) revisite}.
\end{proof}

\begin{Lem}
\label{Prop-clef: existence d'un bon point min successeur}
Let $\eta\in[0,1/2)$. As in Lemma~\ref{Lem: P = QR grand bonne approx}, we set
\[
\eta'' = \frac{(2n-2d+3)\eta+d-3}{n+d-3}.
\]
Suppose either that we have $d=2$, or that we have $d\ge 3$, $\eta''\in [0,1/2)$ and
\begin{equation}
\label{eq Prop: inegalite pour eta}
\frac{1-2\eta''}{1-\eta''} \geq 1-\frac{1}{d-2} + \frac{2d}{n}.
\end{equation}
Then for each large enough $i\geq 0$, there exist $Z\in \bR$ with $\normH{Q_i}\leq Z \leq \normH{Q_i}^{1+\theta_i(1-\eta)}$ and a non-zero $P\in\bZ[X]_{\leq n}$, coprime to $Q_i$, which satisfies $|P(\xi)| < |Q_i(\xi)|$ and
\begin{equation}
\label{eq Prop: system sol premiere avec Q_i}
\normH{P} \leq Z \AND |P(\xi)| \leq Z^{-(2n-d)}.
\end{equation}
\end{Lem}

\begin{proof}
Since $\homega_n(\xi) > 2n-d$, there exists $X_0\geq 0$ such that for each $X \geq X_0$ the system
\begin{equation*}
\normH{P} \leq X \AND |P(\xi)| \leq X^{-(2n-d)}
\end{equation*}
has a non-zero solution $P$ in $\bZ[X]_{\leq n}$. Fix $i\geq 0$ such that $\normH{Q_i}\geq X_0$, and choose a non-zero solution $P\in\bZ[X]_{\leq n}$ of the above system with $X := \normH{Q_i}^{1+\theta_i(1-\eta)}$. For simplicity, write $Q= Q_i$ and $\theta = \theta_i$. We have $|P(\xi)|\leq X^{-(2n-d)} < |Q(\xi)|$ thanks to \eqref{equation: minoration de Q_i(xi) via theta_i}. If $P$ and $Q$ are coprime, then the conclusion holds with $Z=X$. We may therefore assume that $P$ and $Q$ are not coprime. Then $Q$ divides $P$, and assertion~\eqref{item: item 1 Lem P = QR sol } of Lemma~\ref{Lem: P = QR grand bonne approx} implies that $d\geq 3$. Let $A\in\bZ[X]_{\leq d-2}$ and $e\in[1,d-2]$ be the non-constant irreducible polynomial and the integer given by Lemma~\ref{Lem: P = QR grand bonne approx}~\eqref{item: item 2 Lem P = QR sol }. In~particular we have $\deg(A^e)\leq d-2$ and \eqref{eq Prop: P = QR grand bonne approx: H(A^e) revisite} holds. Set $Z:= e^{-2n}\normH{QA^{e}}$, and define~$\nu$ by the relation
\[
Z = \normH{Q}^{1+\theta(1-\nu)}.
\]
By Gelfond's lemma and by definition of $Z$ and $\nu$, we have
\begin{equation*}
\normH{Q}^{\theta(1-\nu)} \asymp \normH{A^e} \gg \normH{Q}^{\theta(1-\eta'')},
\end{equation*}
the last inequality coming from \eqref{eq Prop: P = QR grand bonne approx: H(A^e) revisite}. We deduce that $\nu \leq \eta'' + \GrO\big(1/\log \normH{Q}\big)$. Since $\eta'' < 1/2$ we may assume $i$ large enough so that $\nu < 1/2$. On the other hand, since $QA^e$ divides $P$, by \eqref{eq: Gelfond's lemma}, we have
\[
Z < e^{-n}\normH{QA^e} \leq \normH{P} \leq X = \normH{Q}^{1+\theta(1-\eta)},
\]
hence $\nu\geq \eta$. We now consider a non-zero solution $\tP\in\bZ[X]_{\leq n}$ of the system
\begin{equation}
\label{eq proof: system avec Y}
\normH{\tP} \leq Z \AND |\tP(\xi)| \leq Z^{-(2n-d)}.
\end{equation}
We claim that $\tP$ and $Q$ are coprime. Suppose by contradiction that $Q$ divides $\tP$. There exists $\tR\in\bZ[X]$ such that $\tP = Q\tR$. Write $\tR = A^f\tS$, with $f\in\bN$ and $\tS\in\bZ[X]_{\leq d-2}$ coprime to $A$. By \eqref{eq: Gelfond's lemma} and by definition of $Z$, and since $Q$ and $\tS$ divide $\tP$, we obtain
\[
\normH{Q}\normH{\tS} < e^{n}\normH{\tP} \leq e^nZ = e^{-n} \normH{QA^e} < \normH{Q}\normH{A^e}.
\]
We deduce that $\normH{\tS}\leq \normH{A^e}$. Similarly,
\[
\normH{QA^f} < e^n\normH{\tP} \leq e^nZ = e^{-n} \normH{QA^e}.
\]
Consequently, the polynomial $QA^e$ cannot be a factor of $QA^f$ (by \eqref{eq: Gelfond's lemma} once again). Thus $f\leq e-1$. Since $\normH{\tS}\leq \normH{A^e}$, the last assertion of Lemma~\ref{Lem: P = QR grand bonne approx} yields
\begin{equation}
\label{eq proof : estimation de S(xi)}
|\tS(\xi)| \gg \normH{A^e}^{-(2d-5)}.
\end{equation}
By hypothesis $\nu < 1/2$, and Lemma~\ref{Lem: P = QR grand bonne approx}~\eqref{item: item 1 Lem P = QR sol } applied to the solution $\tP=Q\tR$ of the system \eqref{eq proof: system avec Y} gives the estimate
\begin{equation}
\label{eq proof: minoration omega(S)log H(S)}
|\tR(\xi)| \leq \normH{Q}^{-(n+d-3)(1-\nu')\theta}, \quad \textrm{where } \nu'= \frac{(2n-d)\nu}{n+d-3}.
\end{equation}
We now use \eqref{eq proof : estimation de S(xi)} and $|A^e(\xi)| \gg \normH{A^e}^{-(n+d-3)}$ (coming from \eqref{eq Prop: P = QR grand bonne approx: H(A^e) revisite}) together with $f\leq e-1\leq d-3$. We get the lower bound
\begin{align*}
\log|\tR(\xi)| & = \frac{f}{e}\log|A^e(\xi)|+\log|\tS(\xi)| \\
& \geq -\Big[\Big(1-\frac{1}{e} \Big)(n+d-3) + 2d-5\Big]\log \normH{A^e} + \GrO(1)\\
& \geq -\Big[\Big(1-\frac{1}{d-2} \Big)(n+d-3) + 2d-5\Big]\theta(1-\nu)\log \normH{Q} + \GrO(1),
\end{align*}
the last inequality following from $\normH{A^e} \asymp \normH{Q}^{\theta(1-\nu)}$. Comparing this with \eqref{eq proof: minoration omega(S)log H(S)} and noting that $\nu'\leq 2\nu$, we obtain
\begin{equation*}
\frac{1-2\nu}{1-\nu} \leq \frac{1-\nu'}{1-\nu} \leq 1-\frac{1}{d-2} + \frac{2d-5}{n+d-3} + \GrO\big(1/\log \normH{Q}\big).
\end{equation*}
The function $ \nu \mapsto (1-2\nu)/(1-\nu)$ is decreasing on $[0,1/2]$. Using the estimate $\nu \leq \eta'' + \GrO\big(1/\log \normH{Q}\big)$, we obtain
\begin{equation*}
\frac{1-2\eta''}{1-\eta''} \leq 1-\frac{1}{d-2} + \frac{2d-5}{n+d-3} + \GrO\big(1/\log \normH{Q}\big).
\end{equation*}
Since $(2d-5)/(n+d-3) < 2d/n$, this contradicts our hypothesis \eqref{eq Prop: inegalite pour eta} when $i$ is sufficiently large. So, if $i$ is large enough, then $\tP$ and $Q$ are coprime. Finally, the lower bound $|\tP(\xi)|\leq Z^{-(2n-d)} < |Q(\xi)|$ follows from \eqref{equation: minoration de Q_i(xi) via theta_i} with $\eta$ replaced by $\nu$ (since $\nu < 1/2$), by a similar argument as in the beginning of the proof.
\end{proof}

\begin{proof}[Proof of Proposition~\ref{Prop: estimation omega_n(xi)}]
The condition $d\leq \sqrt[3]{n/4}$ implies that $d\leq 1+n/2$. Define
\begin{equation*}
\eta = \frac{1}{2d+5/2}, \quad \eta'' = \frac{(2n-2d+3)\eta+d-3}{n+d-3} \AND \nu = \frac{1}{d+1}.
\end{equation*}
We claim that the hypotheses of Lemma~\ref{Prop-clef: existence d'un bon point min successeur} are satisfied for this choice of parameters. For $d=2$, this is automatic since $\eta < 1/2$. If $d\geq 3$, a direct computation yields
\begin{equation*}
\eta'' - \nu = \frac{-n+4d^3-11d^2-13d+6}{(4d+5)(n+d-3)(d+1)} < 0,
\end{equation*}
so that $\eta'' < \nu \leq 1/3$. Since $x\mapsto (1-2x)/(1-x)$ is decreasing on $[0,1/2]$, we deduce that $(1-2\eta'')/(1-\eta'') \geq (1-2\nu)/(1-\nu)$. On the other hand, we have
\begin{equation*}
\frac{1-2\nu}{1-\nu} - \Big( 1-\frac{1}{d-2} + \frac{2d}{n}\Big) = \frac{2(n-d^3+2d^2)}{nd(d-2)} \geq 0,
\end{equation*}
hence our claim. Consequently, for each large enough $i$ there exists a non-zero polynomial $P\in\bZ[X]_{\leq n}$ coprime with $Q_i$, satisfying
\[
|P(\xi)| \leq |Q_i(\xi)| < 1 \AND \normH{P}\leq \normH{Q_i}^{1+\theta(1-\eta)}.
\]
Such a polynomial is non-constant, and Lemma~\ref{lem: estimation classique du resultant} yields
\begin{align*}
1 \leq |\Res(Q_i,P)| &\ll \normH{Q_i}^{n-1}\normH{P}^{n}|Q_i(\xi)| + \normH{Q_i}^{n}\normH{P}^{n-1}|P(\xi)| \\
& \ll \normH{Q_i}^{n-1+n(1+\theta(1-\eta))-\omega(Q_i)}.
\end{align*}
As $\normH{Q_i}$ tends to infinity, it follows that
\begin{equation*}
n-1+n(1+\theta(1-\eta))-\omega(Q_i) \geq \GrO\big(1/\log \normH{Q_i} \big).
\end{equation*}
Using the definition \eqref{eq: rappel theta_i} of $\theta_i$, this can be rewritten as
\begin{equation*}
(n\eta-2d+3)\omega(Q_i) \leq 2\eta n^2-(3d+\eta d-5)n+2d-3 + \GrO\big(1/\log \normH{Q_i} \big).
\end{equation*}
The hypothesis $d\leq \sqrt[3]{n/4}$ implies $n\eta-2d+3 > 0$. Thus, after simplification
\begin{align*}
\omega(Q_i) + \GrO\Big(1/\log \normH{Q_i} \Big) &\leq \frac{2\eta n^2-(3d+\eta d-5)n+2d-3}{n\eta-2d+3} \\
& \; = 2n + \frac{n(d-1-\eta d)+2d-3}{n\eta-2d+3} = 2n + P(n,d),
\end{align*}
where $P(n,d)$ is as in the statement of Proposition~\ref{Prop: estimation omega_n(xi)} (and $\eta=1/(2d+5/2)$). We~conclude that
\[
\omega_n(\xi) = \limsup_{i\to\infty} \omega(Q_i) \leq 2n+P(n,d).
\]
Set $Q(n,d) = (2n-8d^2+2d+15)(P(n,d)-2d^2)$. A direct computation yields
\begin{equation*}
Q(n,d) = -n(d+5)+16d^4-4d^3-22d^2-2d-15.
\end{equation*}
If $\displaystyle d\leq \sqrt[3]{n/16}$ we have $16d^4 \leq nd$, and therefore $Q(n,d)\leq 0$. We obtain $P(n,d)\leq 2d^2$, and consequently $\homega_n(\xi)\leq 2n+2d^2$. It remains to show that in the case $n\geq 17$ and $d= \lceil \rho \rceil$ with
\[
\rho = \sqrt[3]{n/16},
\]
we still have $Q(n,d)\leq 0$. If $17 \leq n \leq 128$, or equivalently if $1 < \rho \leq 2$, then we have $d=2$ and $Q(n,2) = -7n+117 \leq 0$. The same reasoning leads to $Q(n,d) \leq 0$ for $2 < \rho \leq 3$ and $3 < \rho \leq 4$. We now suppose that $\rho > 4$. Writing $d = \rho + t$, with $t\in [0,1]$, and using the fact that $16\rho^3 = n$, we find
\begin{align*}
Q(n,d) \leq -n(d+5)+16d^4 & = - 16\rho^3(\rho+t+5) + 16 \big(\rho^4 + 4t\rho^3 + 6t^2\rho^2 + 4t^3\rho +t^4 \big) \\
& = 16\rho^3(3t-5) + 16 \big(6t^2\rho^2 + 4t^3\rho +t^4 \big) \leq 16R(\rho),
\end{align*}
where $R(x) = -2x^3+6x^2+4x+1$. As the coefficients of $R(x+4)$ are all negative, we~have $R(x)\leq 0$ for each $x\geq 4$. In particular, $R(\rho)\leq 0$, and once again we~obtain $Q(n,d)\leq 0$.
\end{proof}

Note that the upper bound $2n+P(n,d)$ is not optimal in Proposition~\ref{Prop: estimation omega_n(xi)} (and could be slightly improved by choosing the parameter $\eta$ closer to $1/(2d)$).

\section{Proof of the main theorem}
\label{section: proof of main thm}

In this last section we prove our main Theorem~\ref{Thm : main} in the following stronger form.

\begin{Thm}
\label{Thm : main avec meilleure cst}
Let $\ee = 0.3748\cdots$ be the unique (positive) solution of the equation $(1+x)e^x = 2$ and set $a = \big(2\ee(2-e^\ee)/9\big)^{1/3} = 0.3567 \cdots$. There exists an explicit constant $C>0$ such that, for each $n\geq 1$ and any transcendental real number $\xi\in\bR$, we have
\begin{equation*}
\homega_n(\xi) \leq 2n- a n^{1/3} + C.
\end{equation*}
\end{Thm}

Since $1/3 < a$, it implies Theorem~\ref{Thm : main}. We first establish a preliminary result which uses the following notation. Let $n,d$ be integers with $2\leq d \leq \sqrt{n/4}$. In particular $d\leq 1+n/2$. We define
\[
\omega(d,n) := 2n+P(n,d), \quad \textrm{where } P(n,d) = \frac{n(4d^2-d-5)+8d^2-2d-15}{2n-8d^2+2d+15},
\]
as well as
\[
\tau(d,n) = \frac{(2n-d)(n-2d+3)}{\omega(d,n)\big(\omega(d,n)-n-d+3\big)} \AND \mu(d,n):= \frac{n}{(2n-d)\tau}.
\]
Let $\big(\tau_i(d,n)\big)_{0\leq i\leq n/2}$ be the sequence associated to $\tau=\tau(d,n)\in(0,1)$ by Definition~\ref{Def: tau_j}.

\begin{Thm}
\label{Thm: main theorem under condition}
Let $n,d, j$ be non-negative integers with $2\leq d \leq \sqrt{n/4}$ and $1\leq j\leq n/2$. Suppose that
\begin{equation}
\label{eq: condition sur j}
\tau_k(d,n) > \mu(d,n) \quad \textrm{for $k=0,\dots,j$}.
\end{equation}
Then for any transcendental real number $\xi$ we have
\begin{equation}
\label{eq proof: resultat main a montrer}
\homega_n(\xi)\leq 2n- \min\big\{d,d_j\big\}, \quad \textrm{where } d_j= 2j-1-\frac{j-1}{\tau_j(d,n)}.
\end{equation}
\end{Thm}

\begin{proof}
Fix a transcendental real number $\xi$. If $\homega_n(\xi)\leq 2n-d$, then \eqref{eq proof: resultat main a montrer} holds. We now assume that $\homega_n(\xi) > 2n-d$, and we choose a real number $\homega$ such that
\[
\homega_n(\xi) > \homega > 2n-d.
\]
Let $(P_i)_{i\geq 0}$ denote a sequence of minimal polynomials associated to $n$ and $\xi$ as in Section~\ref{section: minimal pol}. We denote by $(Q_i)_{i\geq 0}$ the sequence of irreducible factors given by Proposition~\ref{prop: existence Qi}, and denote by
\[
m:=m_n(\xi)
\]
the dimension of the spaces $\Vect[\bR]{Q_i,Q_{i+1},\dots}$ for each large enough $i$ (as in Definition~\ref{Def m_n}). Proposition~\ref{Prop: estimation omega_n(xi)} yields $\omega_n(\xi)\leq \omega(d,n)$, and by Proposition~\ref{prop: existence Qi}~\eqref{enum: H(Q_i+1) controle par H(Q_i)}, we get, for each large enough $i$,
\[
\normH{Q_{i+1}}^{\tau(d,n)} \leq \normH{Q_i}.
\]
For simplicity, we write $\tau=\tau(d,n)$ and $\tau_k=\tau_k(d,n)$ for each $k\in\bN$ with $k\leq n/2$. If $m-2 \le j\le n/2$, then Corollary~\ref{Cor: tau_m_2 < mu} yields $\tau_{m-2} \leq \mu(d,n)$, which contradicts the hypothesis \eqref{eq: condition sur j}. Hence, we must have $j<m-2$. Let $i\geq 0$. Set $Q = Q_{\sigma_j(i)}$. If $i$ is large enough, there exists a non-zero $P\in\bZ[X]_{\leq n}$ such that
\begin{equation*}
\normH{P} \leq e^{-n} \normH{Q}=: X \AND |P(\xi)| \leq X^{-\homega}.
\end{equation*}
By \eqref{eq: Gelfond's lemma} the (irreducible) polynomial $Q$ does not divide $P$, they are thus coprime. Lemma~\ref{lem : dim V_n+j(P,Q)} implies that $\dim V_{2n-j}(P,Q)\geq 2n-2j+2$. Choose a linearly independent subset
\[
\cU_j:=\{U_1,\dots,U_{2n-2j+2}\} \subset \cB_{2n-j}(P,Q)
\]
of cardinality $2n-2j+2$. According to \eqref{eq: V_2n-j(A_j[i]) = tout l'espace}, we have $V_{2n-j}(A_j[i]) = \bR[X]_{\leq 2n-j}$. So~there exists
\[
\cV_j:=\{V_1,\dots,V_{j-1}\} \subset \cB_{2n-j}(A_j[i]) = \cB_{2n-j}(Q_{\sigma_j(i)},\dots, Q_i)
\]
such that
\[
\Vect[\bR]{\cU_j} \oplus \Vect[\bR]{\cV_j} = \bR[X]_{\leq 2n-j}.
\]
Then, identifying $\bR[X]_{\leq 2n-j}$ with $\bR^{2n-j+1}$ via \eqref{eq: identification poly avec points de R^m}, we form the determinant
\begin{equation}
\label{eq proof: det à considérer}
1 \leq \big|\det(U_1,\dots,U_{2n-2j+2},V_1,\dots,V_{j-1}) \big|.
\end{equation}
For $k=1,\dots,2n-2j+2$, we have
\begin{equation*}
\normH{U_k} \ll \normH{Q} \AND |U_k(\xi)| \ll \normH{Q}^{-\homega}.
\end{equation*}
On the other hand, for $k=1,\dots,j-1$, we have by Equation \eqref{eq Prop: minoration de Y_j} from Proposition~\ref{Prop: prop resumee avec tau_j}
\begin{equation*}
\normH{Q} \ll \normH{V_k} \ll \normH{Q_i} \ll \normH{Q}^{1/\tau_j} \AND |V_k(\xi)| \ll \normH{V_k}^{-\homega} \ll \normH{Q}^{-\homega}.
\end{equation*}
For $i=2,\dots,2n-j+1$, we add to the first row of the determinant \eqref{eq proof: det à considérer} the $i$-th row multiplied by $\xi^{i-1}$. This first row now becomes
\[
\big(U_1(\xi),\dots,U_{2n-2j+2}(\xi),V_1(\xi),\dots,V_{j-1}(\xi)).
\]
By the above, the absolute value of each of its elements is $\ll \normH{Q}^{-\homega}$. By expanding the determinant, we obtain
\begin{equation*}
1 \ll \normH{Q}^{2n-2j+1}\normH{Q_i}^{j-1}\normH{Q}^{-\homega} \ll \normH{Q}^{2n-2j+1+(j-1)/\tau_j-\homega}.
\end{equation*}
By letting $i$ tend to infinity, we deduce that
\begin{equation*}
\homega \leq 2n-2j+1+(j-1)\tau_j^{-1} = 2n-d_j.
\end{equation*}
Since $\homega$ may be chosen arbitrarily close to $\homega_n(\xi)$, we finally get \eqref{eq proof: resultat main a montrer}.
\end{proof}

In view of \eqref{eq proof: resultat main a montrer}, the idea is now to choose $d$ and $j$ so that $d$ is maximal and $d \thickapprox d_j$. The next two results aim at simplifying condition~\eqref{eq: condition sur j} of Theorem~\ref{Thm: main theorem under condition}. The second one also provides a simple lower bound for the exponent $\tau_j$.

\begin{Lem}
\label{lem: condition thm clef seulement pour k=j}
Let $n,d, j$ be non-negative integers with $2\leq d \leq \sqrt{n/4}$ and $ 1\leq j\leq n/2$. Suppose that $j$ satisfies
\begin{equation}
\label{eq lem: condition sur k=j seulement pour tau_j}
\frac{(2n-d)\tau(d,n)^2}{(n-2j)\tau(d,n)+n-j+1} \leq 1 \AND \tau_j(d,n) \geq 0.
\end{equation}
Then, the sequence $\big(\tau_k(d,n)\big)_{0\leq k \leq j}$ is (strictly) decreasing. In particular, condition~\eqref{eq: condition sur j} is fulfilled if moreover
\begin{equation*}
\tau_j(d,n) > \mu(d,n).
\end{equation*}
\end{Lem}

\begin{proof}
Let $\alpha_1 \leq \cdots \leq \alpha_j$ be as in Definition~\ref{Def: tau_j}. Condition \eqref{eq lem: condition sur k=j seulement pour tau_j} is equivalent to $\alpha_j\leq 1$ and $\tau_j(d,n) \geq 0$. By definition, we have
\[
\tau_{k-1}(d,n) = \alpha_k^{-1}\tau_k(d,n) + \frac{2k-1}{2n-d} \qquad \textrm{(for $k=1,\dots,j$).}
\]
Since $\alpha_k^{-1}\geq \alpha_j^{-1}\geq 1$, this yields $\tau_{k-1}(d,n) > \tau_k(d,n)$. This proves the first assertion of our lemma. The second one follows easily.
\end{proof}

\begin{Lem}
\label{lem: minorer tau_j}
Let $n,d, j$ be non-negative integers with $2\leq d \leq \sqrt{n/4}$ and $ 1\leq j\leq n/2$. Define
\[
\alpha = \alpha(d,n) := \frac{(2n-d)\tau(d,n)^2}{(n-2)\tau(d,n)+n},
\]
and suppose that
\begin{equation}
\label{eq lem: condition pour minorer tau_j}
\alpha^{j} > \frac{j(2j-1)\tau(d,n)}{(n-2)\tau(d,n)+n} = \frac{j(2j-1)\alpha}{(2n-d)\tau(d,n)}.
\end{equation}
Then, $\alpha\in(0,1)$ and for $k=0,\dots,j$, we have
\begin{equation*}
\tau_k(d,n) \geq \alpha^j\tau(d,n) - \frac{j(2j-1)\tau(d,n)^2}{(n-2)\tau(d,n)+n} > 0.
\end{equation*}
\end{Lem}

\begin{proof}
We have $\alpha\in(0,1)$ since $\tau(d,n)< 1$ and $d\geq 2$. For simplicity, we write $\tau = \tau(d,n)$. Let $(\sigma_k)_{k\geq 0}$ be the sequence defined by $\sigma_0= \tau$, and
\begin{equation*}
\sigma_k = \alpha\Bigl(\sigma_{k-1} - \frac{2k-1}{2n-d} \Bigr) \qquad \textrm{for } k\geq 1.
\end{equation*}
Using \eqref{eq lem: condition pour minorer tau_j}, we find
\begin{equation}
\label{eq proof: minoration sigma_j}
\frac{\sigma_j}{\alpha^j} = \frac{\sigma_{j-1}}{\alpha^{j-1}} - \frac{2j-1}{(2n-d)\alpha^{j-1}} = \sigma_0 -\frac{1}{2n-d}\sum_{k=1}^{j} \frac{2k-1}{\alpha^{k-1}} \geq \tau -\frac{j(2j-1)}{(2n-d)\alpha^{j-1}} > 0.
\end{equation}
In particular $\sigma_j\geq 0$. Since $\sigma_{k-1} \geq \alpha^{-1}\sigma_k$, by induction, we get $\sigma_j < \sigma_{j-1} < \cdots < \sigma_0$. Moreover, $\alpha = \alpha_1 \leq \alpha_k$, for each $k\in\bN$ with $1\leq k \leq n/2$, where $\alpha_k$ is as in Definition~\ref{Def: tau_j}. Another quick induction yields $\sigma_k\leq \tau_k$ for $k=0,\dots,j$. We conclude by combining $\sigma_j\leq \sigma_k\leq \tau_k$ with \eqref{eq proof: minoration sigma_j}.
\end{proof}

\begin{proof}[Proof of Theorem~\ref{Thm : main avec meilleure cst}]
Define a function $f:[0,\infty)\to \bR$ by $f(x) = x(2-e^x)$. Let $\ee = 0.3748\cdots$ be the unique solution of the equation $(1+x)e^x = 2$. It is the abscissa of the maximum of $f$. Set
\begin{equation*}
a = \sqrt[3]{\frac{2\ee(2-e^\ee)}{9}} = 0.3567\cdots
\end{equation*}
Let $n\geq 1$ and define $d=d(n)$ and $j(n)$ by
\begin{equation*}
d(n) = \lceil a n^{1/3} \rceil \AND j:=j(n):= \bigg\lceil \frac{2\ee n}{9d^2} \bigg\rceil.
\end{equation*}
We suppose $n \geq 30$ so that $2 \leq d \leq 1+n/2$ and $1\leq j\leq n/2$. Since $d^4/n^2 \asymp d/n = \GrO(n^{-2/3})$, we find $\omega(d,n) = 2n+2d^2 + \GrO(d)$, and then
\begin{equation*}
\tau(d,n) = 1-\frac{3d^2}{n} +\GrO(\sfrac{1}{n^{2/3}}) \AND \alpha(d,n) = 1-\frac{9d^2}{2n} +\GrO(\sfrac{1}{n^{2/3}}),
\end{equation*}
(where $\alpha(d,n)$ is defined in Lemma~\ref{lem: minorer tau_j}). In particular, by choice of $j$, we have
\begin{equation}
\begin{aligned}
\label{eq proof: minoration alpha^j}
\alpha(d,n)^j = \exp\big(j\log(\alpha(d,n))\big) & = \exp\Bigl(-\frac{9jd^2}{2n} +\GrO(\sfrac{1}{n^{1/3}}) \Bigr)\\
& = e^{-\ee} + \GrO(\sfrac{1}{n^{1/3}}).
\end{aligned}
\end{equation}
Since
\[
\frac{j(2j-1)\tau(d,n)}{(n-2)\tau(d,n)+n} = \GrO(\sfrac{1}{n^{1/3}}),
\]
there exists $N_1\geq 30$ such that condition \eqref{eq lem: condition pour minorer tau_j} of Lemma~\ref{lem: minorer tau_j} is fulfilled for each $n\geq N_1$. Thus, for $k=0,\dots,j$, we have
\begin{equation*}
\tau_k(d,n) \geq \alpha(d,n)^j\tau(d,n) - \frac{j(2j-1)\tau(d,n)^2}{2n-2} = e^{-\ee} + \GrO(\sfrac{1}{n^{1/3}}).
\end{equation*}
In particular $d_j = 2j-1-(j-1)/\tau_j(d,n)$ satisfies
\begin{equation*}
d_j \geq j\big(2- e^{\ee}\big) + \GrO(1) = \frac{2\ee(2- e^{\ee}) n}{9d^2} + \GrO(1) = \frac{a^3 n}{d^2} + \GrO(1) = d + \GrO(1).
\end{equation*}
On the other hand, we have
\begin{equation*}
\mu(d,n) = \frac{n}{(2n-d)\tau} = \frac{1}{2}+\GrO(\sfrac{1}{n^{1/3}}).
\end{equation*}
Since $e^{-\ee} > 1/2$, by \eqref{eq proof: minoration alpha^j} there exists $N_2\geq N_1$ such that condition \eqref{eq: condition sur j} of Theorem~\ref{Thm: main theorem under condition} is fulfilled for each $n\geq N_2$. We conclude that for any $n\geq N_2$ and any transcendental real number $\xi$, we have
\begin{equation*}
\homega_n(\xi) \leq 2n-\min\{d,d_i\} = 2n - d + \GrO(1).\qedhere
\end{equation*}
\end{proof}

\appendix

\section*{Appendix. Twisted heights}
\refstepcounter{section}

The purpose of this appendix is to give another interpretation of the quantity $\HHstar(V)$ defined in Section~\ref{def petitesse HHstar(V)}. Our first approach was actually to work with the heights $\HH_T$ defined below. We are thankful to Damien Roy for pointed out the link with Hodge's duality.

Fix $A\in \GL(\bR^{m+1})$ and let $V$ be a $k$-dimensional subspace of $\bR^{m+1}$ defined over~$\bQ$. Its (twisted) height $\HH_A(V)$ is defined as
the covolume of the lattice $A(V\cap\bZ^{m+1})$ inside the subspace $A(V)$ (with the convention that $\HH_A(V) = 1$ if $V=\{0\}$). Explicitly, we~have
\begin{equation}
\label{eq: def hauteur tordue}
\HH_A(V) := \norm{A\bx_1\wedge \cdots \wedge A\bx_k },
\end{equation}
where $(\bx_1,\dots,\bx_k)$ is any $\bZ$-basis of the lattice $V\cap \bZ^{m+1}$. Then Schmidt's inequality generalizes as follows
\begin{equation}
\label{eq : inegalite de Schmidt}
\HH_A(U+V) \HH_A(U\cap V) \leq \HH_A(U) \HH_A(V)
\end{equation}
for any subspaces $U,V$ of $\bR^{m+1}$ defined over $\bQ$. The proof is the same as for rational subspaces (see \cite[Ch.\,I, Lem.\, 8A]{Schmidt} and \cite[\S 5]{marnat2018optimal}). Similarly to Marnat and Moshchevitin \cite[\S 5]{marnat2018optimal}, we consider twisted heights of the following form. Let $T>1$ be a parameter. We define the matrix $A_{m,T}\in\GL(\bR^{m+1})$ as
\[
A_{m,T} = \begin{pmatrix}
T^{m} & 0 & \dots & 0 \\
0 & T^{-1} & & \\
\vdots & & \ddots & \\
0 & & & T^{-1}
\end{pmatrix}
\begin{pmatrix}
1 & \xi & \cdots & & \xi^{m} \\
0 & 1 & 0 & \cdots & 0 \\
\vdots & & \ddots & & \vdots \\
& & & & \\
0 & \cdots & & 0 & 1
\end{pmatrix},
\]
so that for each polynomial $P= a_0+\cdots + a_{m}X^{m} \in \bZ[X]_{\leq m}$ (identified to a point of $\bR^{m+1}$ via \eqref{eq: identification poly avec points de R^m}), we have
\begin{equation}
\label{eq: expression de Ax pour la hauteur T}
A_{m,T}\begin{pmatrix}
a_0 \\
\vdots \\
a_{m}
\end{pmatrix}=
\begin{pmatrix}
T^{m}P(\xi) \\
T^{-1}a_1 \\
\vdots \\
T^{-1} a_m
\end{pmatrix}.
\end{equation}
We denote by $\HH_{m,T}$ (or simply $\HH_T$ if there is no ambiguity about the integer $m$) the twisted height $\HH_{A}$ associated to the matrix $A=A_{m,T}$. Note that
\[
\HH_{T}\big(\bR[X]_{\leq m}\big) = \HH_{T}(\bR^{m+1}) = \det(A) = 1.
\]

\begin{Def}
Let $V$ be a subspace of $\bR[X]_{\leq m}$ defined over $\bQ$. We set
\begin{equation*}
\HHH(V) = \lim_{T\to+\infty} T^{-\codim(V)}\HH_{m,T}(V),
\end{equation*}
where $\codim(V) = m+1-\dim(V)$ denotes the codimension of the space $V$ inside $\bR[X]_{\leq m}$. In particular, $\HHH(\bR[X]_{\leq n}) = 1$, and for any primitive polynomial $P\in \bZ[X]_{\leq m}$, we have
\[
\HHH\big(\Vect[\bR]{P}\big) = |P(\xi)| = \HHstar\big(\Vect[\bR]{P}\big).
\]
\end{Def}

Our goal is now to prove that for any non-zero subspace $V\subset\bR[X]_{\leq m} \simeq \bR^{m+1}$ defined over $\bQ$, we have
\[
\HHH(V) \asymp \HHstar(V),
\]
where $\HHstar$ is as in Definition~\ref{Def: def HHstar} (and the implicit constant depends on $m$ and $\xi$ only). First, note that since $\dim(U+V) + \dim(U\cap V) = \dim U + \dim V$ for any subspaces $U,V$ of $\bR[X]_{\leq m}$, we deduce from \eqref{eq : inegalite de Schmidt} (with $A=A_{m,T}$) the following version of Schmidt's inequality, which is the analog of Proposition~\ref{Prop: Schmidt inegalite pour HHstar}
\begin{equation}
\label{eq : inegalite de Schmidt new}
\HHH(U+V) \HHH(U\cap V) \leq \HHH(U) \HHH(V),
\end{equation}
valid for any $U,V$ of $\bR[X]_{\leq m}$ defined over $\bQ$.

\begin{Prop}
\label{Prop: annexe, formule HHH(V)}
Let $V$ be a $k$-dimensional subspace of $\bR^{m+1}$ defined over $\bQ$, with $1 \leq k\leq m+1$, and set $\Xi_m=(1,\xi,\dots,\xi^m)$. We have
\begin{equation}
\label{Prop annexe: equivalence estimate avec hauteur tordue}
\HHstar(V) \ll \HHH(V) \leq \HHstar(V),
\end{equation}
where the implicit constant depends on $\xi$ and $m$ only. Moreover, for any $\bZ$-basis $(\bx_1,\dots,\bx_k)$ of $V\cap\bZ^{m+1}$, we have
\begin{equation}
\label{eq prop: formule explicite HHH(V)}
\HHH(V) = \bnorm{\sum_{i=1}^{k} (-1)^{k-i}\psc{\Xi_m}{\bx_i}\,\bx_1^+\wedge \cdots \wedge\widehat{\bx_i^+}\wedge \cdots \wedge\bx_k^+},
\end{equation}
where $\bx_i^+\in\bZ^{m}$ denotes the point $\bx_i$ deprived of its first coordinate.
\end{Prop}

Before to prove this result, we introduce some notation that we will need in the proof. Given two positive integers $p$ and $q$, we define $\cI(p,q)$ as the set of $p$-tuples $(i_1,\dots,i_p)$ of integers with $1\leq i_1 < \cdots < i_p \leq q$. Let $\be = (\be_1,\dots,\be_q)$ be the canonical basis of $\bR^q$. For any $I\in\cI(p,q)$ as above, set $\be_{I} = \be_{i_1}\wedge \cdots \wedge \be_{i_p} \in \sbw{p}{\bR^q}$. For any $\bX\in\sbw{p}{\bR^q}$, we call \emph{$I$-coordinate of $\bX$} its $\be_I$-coordinate in the basis $(\be_J)_{J\in\cI(p,q)}$. For any $\bx_1,\dots,\bx_p\in\bR^q$, we denote by $M(\bx_1,\dots,\bx_p)$ the $q\times p$ matrix whose columns are $\bx_1,\dots, \bx_p$ written in the basis $\be$, and by $\cD_{I}(\bx_1,\dots,\bx_p)$ the minor formed by the rows of $M(\bx_1,\dots,\bx_p)$ of index $i$ in $I$. Then, writing $\bX = \bx_1\wedge \cdots \wedge\bx_k$, we have the classical formulas
\begin{equation}
\label{eq: formule norm multi-vecteur}
\bX = \sum_{I\in\cI(p,q)}\cD_{I}(\bx_1,\dots,\bx_p) \be_{I} \AND \norm{\bX}^2 = \sum_{I\in\cI(p,q)}\cD_{I}(\bx_1,\dots,\bx_p)^2.
\end{equation}
Therefore, for each $I\in\cI(p,q)$, the $I$-coordinate of $\bX$ is $\cD_{I}(\bx_1,\dots,\bx_p)$.

\begin{proof}[Proof of Proposition~\ref{Prop: annexe, formule HHH(V)}]
Fix $T\geq 1$ and for $i=1,\dots,k$, set
\begin{align*}
\bZZ &= \sum_{i=1}^kp_i\bx_1\wedge \cdots \wedge\widehat{\bx_i}\wedge \cdots \wedge\bx_k, \quad \textrm{where } p_i = (-1)^{i+1}\psc{\Xi_m}{\bx_i},\\
\bY & = \lim_{T\to+\infty} T^{-m+k-1}\by_1(T)\wedge\cdots\wedge \by_k(T), \quad \textrm{where } \by_i = \by_i(T) = A_{m,T}(\bx_i) \in \bR^{m+1}.
\end{align*}
By \eqref{eq: HHstar(V) formule explicite} we have
\begin{equation*}
\HHstar(V) = \norm{\bZZ} \AND \HHH(V) = \norm{\bY}.
\end{equation*}
We prove the following properties. For $i=1,\dots,k$ we set $\bz_i = (\psc{\Xi_m}{\bx_i}, \bx_i)\in \bR^{m+2}$.
\begin{enumerate}
\item \label{enum preuve annexe: item 1} For each $J=(1,j_2,\dots,j_k)\in\cI(k,m+2)$, the $J$-coordinate of $\bz_1\wedge\cdots \wedge\bz_k$ is equal to the $K$-coordinate of $\bZZ$, where $K=(j_2-1,\dots,j_k-1)\in\cI(k-1,m+1)$.
\end{enumerate}
Fix $I = (i_1,\dots,i_k)\in\cI(k,m+1)$.
\begin{enumerate}
\setcounter{enumi}{1}
\item \label{enum preuve annexe: item 2} If $i_1 \geq 2$, then the $I$-coordinate of $\bY$ is equal to $0$.
\item \label{enum preuve annexe: item 3} If $i_1=1$, then the $I$-coordinate of $\bY$ is equal to the $J$-coordinate of $\bz_1\wedge\cdots\wedge\bz_k$, where $J = (1,i_2+1,\dots,i_k+1)$. It is also equal to the $K$-coordinate of $\bZZ$, where $K=(i_2,\dots,i_k)$.
\end{enumerate}
To prove the first assertion, it suffices to expand the determinant $\cD_{J}(\bz_1,\dots,\bz_k)$ along its first row. Let $I = (i_1,\dots,i_k)\in\cI(k,m+1)$. Suppose first that $i_1\neq 1$. Then, by~Hadamard's inequality, the $I$-coordinate of $\by_1(T)\wedge\cdots \wedge\by_k(T)$ satisfies
\[
|D_{I}(\by_1,\dots,\by_k)| \ll \prod_{j=1}^{k} T^{-1}\normH{\bx_j} = \GrO(T^{-k}),
\]
and we deduce that the $I$-coordinate of $\bY$ is equal to $0$, which proves assertion~\eqref{enum preuve annexe: item 2}. Suppose now that $i_1=1$ and set $J = (1,i_2+1,\dots,i_k+1)$. Then
\begin{equation*}
D_{I}(\by_1,\dots,\by_k) = T^{m+1-k}\cD_J(\bz_1,\dots,\bz_k),
\end{equation*}
hence the first part of~\eqref{enum preuve annexe: item 3}. The second part is obtained by combining the above with assertion~\eqref{enum preuve annexe: item 1}.

We deduce from the last two assertions that all the non-zero coordinates of $\bY$ are coordinates of $\bZZ$, thus $\norm{\bY} \leq \norm{\bZZ}$, which proves the second inequality in \eqref{Prop annexe: equivalence estimate avec hauteur tordue}. For the first estimate, we need to estimate the $K$-coordinates of $\bZZ$ with $K\in\cI(k-1,m+1)$ of the form $(1,i_2,\dots,i_{k-1})$. According to assertion~\eqref{enum preuve annexe: item 1}, they are exactly the determinants $\cD_J(\bz_1,\dots,\bz_k)$ with $J = (1,2,j_3,\dots,j_k)$ in $\cI(k,m+2)$.

Fix a $J\in \cI(k,m+2)$ as above. The second row of the matrix $M(\bz_1,\dots,\bz_k)$ is a linear combination of the remaining rows (with coefficients in absolute value between~$1$ and~$|\xi|^m$). We deduce that $\cD_J(\bz_1,\dots,\bz_k)$ can be written as a linear combination of $\cD_{J'}(\bz_1,\dots,\bz_k)$, where $J'$ belong to the subset of $\cI(k,m+2)$ consisting in the $k$-tuples whose second element is $\geq 3$. By assertion~\eqref{enum preuve annexe: item 3}, they are all coordinates of $\bY$, hence $|\cD_J(\bz_1,\dots,\bz_k)|\ll \norm{\bY}$. We conclude that $\norm{\bZZ}\ll \norm{\bY}$.

Finally, fix $(i_2,\dots,i_k)\in\cI(k-1,m)$ and set $K = (i_2+1,\dots,i_k+1)$. By definition of $\bZZ$, the $K$-coordinate of $\bZZ$ is equal to
\begin{equation*}
\sum_{i=1}^{k}p_i \cD_I(\bx_1,\cdots,\widehat{\bx_i}, \cdots,\bx_k) = \sum_{i=1}^{k}p_i \cD_J(\bx_1^+,\cdots,\widehat{\bx_i^+}, \cdots,\bx_k^+).
\end{equation*}
By assertion~\eqref{enum preuve annexe: item 3}, this is also the $(1,K)$-coordinate of $\bY$. So, the set of non-zero coordinates of $\bY$ is exactly equal to the set of non-zero coordinates of the point
\[
\sum_{i=1}^{k}p_i\,\bx_1^+\wedge \cdots \wedge\widehat{\bx_i^+}\wedge \cdots \wedge\bx_k^+.
\]
Equation \eqref{eq prop: formule explicite HHH(V)} follows from the second identity of \eqref{eq: formule norm multi-vecteur}.
\end{proof}

\backmatter
\bibliographystyle{jepplain+eid}
\bibliography{poels}
\end{document} 